\documentclass[12pt,a4paper]{report}


% --- Langue / encodage ---
\usepackage[T1]{fontenc}
\usepackage[utf8]{inputenc} % OK avec pdflatex
\usepackage[french]{babel}

\addto\captionsfrench{%
	\renewcommand{\figurename}{Carte}%
	\renewcommand{\listfigurename}{Liste des cartes}%
}

% --- Mise en page minimale ---
\usepackage[a4paper,margin=2.5cm]{geometry}
\usepackage{microtype}
\usepackage{lmodern}
\sloppy

% --- Autorise les underscores dans le texte ---
\usepackage[strings]{underscore}

% --- Figures / tableaux ---
\usepackage{graphicx}
\usepackage{booktabs}
\usepackage{longtable}
\usepackage{float}


% --- Hyperliens (ne charger QU'UNE fois) ---
\usepackage[hypertexnames=false]{hyperref}
\hypersetup{
	colorlinks=true,
	linkcolor=blue,
	urlcolor=blue,
	citecolor=blue
}

% --- Bibliographie en notes complètes ---
\usepackage{csquotes}
\usepackage[
	backend=biber,
	style=verbose-trad2,
	citetracker=true
	]{biblatex}
\addbibresource{bibliographie.bib}
\usepackage{xurl} % coupe mieux les URL
\usepackage{hyperref} 


% --- Index multiples ---
\usepackage{imakeidx}

\makeindex[name=nom,     title={Index des noms}]
\makeindex[name=lieux,   title={Index des lieux}]
\makeindex[name=auteur,  title={Index des auteurs}]
\makeindex[name=tab,     title={Index des tableaux}, columns=1]
\makeindex[name=carte,   title={Index des cartes et illustrations}, columns=1]
\makeindex[name=image,   title={Index des images}, columns=1]

% --- Commandes d’indexation ---
\newcommand{\idxnom}[1]{\index[nom]{#1}}
\newcommand{\idxlieu}[1]{\index[lieux]{#1}}
\newcommand{\idxauteur}[1]{\index[auteur]{#1}}
\newcommand{\idxtab}[1]{\index[tab]{#1}}
\newcommand{\idxcarte}[1]{\index[carte]{#1}}
\newcommand{\idximage}[1]{\index[image]{#1}}

\renewcommand{\figurename}{Fig.}

% --- Helpers pour légendes + index ---

% Table : caption + label + index "tab"
\newcommand{\tabcaptionindex}[2]{%
	\caption{#2}\label{#1}\idxtab{#2|hyperpage}%
	\idxtab{Tableau \thetable\space -- #2|hyperpage}%
}

%old Figure : caption + label + index selon type ("carte" ou "image")
% Usage : \figcaptionindex{carte}{fig:...}{Texte}
\newcommand{\figcaptionindex}[3]{%
	\caption{#3}\label{#2}%
	\def\tempa{#1}\def\tempb{Carte}%
	\ifx\tempa\tempb
	\idxcarte{#3|hyperpage}%
	\else
	\def\tempb{image}%
	\ifx\tempa\tempb
	\idximage{#3|hyperpage}%
	\fi
	\fi
}



% --- Macro figure simple (ta macro) ---
% IMPORTANT : l'argument optionnel doit être collé : \cartefigure[0.5\textwidth]{...}{...}{...}
\newcommand{\cartefigure}[4][0.7\textwidth]{%
	\begin{figure}[H]
		\centering
		\includegraphics[width=#1]{#2}
		\caption{#4}
		\label{#3}
		% Index : "Carte 3.1 — Titre"
		\idxcarte{Fig.~\thefigure\space--\space #4|hyperpage}%
	\end{figure}%
}



%%%%%%%%%%%%%%%%%%%%%%%%%%%%%%%%%%%%%%%%%%%%%%%%%%%%%%%%%%%%%%%
\begin{document}
	
	\title{Propriété des alpages à Saint-Maxime de Beaufort au 18e siècle}
	\author{Nicolas MORANE}
	\date{\today}
	\maketitle
	
	\tableofcontents
	
	
\chapter*{Table des Abréviations}
\addcontentsline{toc}{chapter}{Table des Abréviations}


ADS	Archives Départementales de Savoie

ADHS 	Archives Départementales de Haute-Savoie

AHSMA : Archives Historique de l'abbaye de Saint Maurice d'Agaune
	
\chapter*{Remerciements}
\addcontentsline{toc}{chapter}{Remerciements}
	
	
	C'est vers le Beaufortain et vers ses habitants que je voudrais adresser en premier lieu mes remerciements. 
	Je remercie mon épouse, étrangère à la montagne, d'avoir accepté et aimé à son tour cette vallée unique.
	Évidemment je remercie les conseils et enseignements de mes professeurs de la Sorbonne, et en particulier de m'avoir accepté dans ce cursus de Master, un étudiant plus tout jeune.
	Je remercie aussi, Philippe Matignon, descendant du vieil allobroge, le sénateur Pierre Blanc, pour m'avoir montré et permis d'utiliser ici, cette très grande carte, peinte sur toile qu'il a sortie d'une armoire de sa maison de Beaufort. 
	
	
\chapter*{Introduction}
\addcontentsline{toc}{chapter}{Introduction}


Dans la plus grande maison de Beaufort (autrefois Saint-Maxime de Beaufort) se trouve une grande toile enroulée sur elle-même, lorsqu'on la déroule (1,40m sur 2m), on aperçoit des motifs peint à la main en couleur, représentant des parcelles et tout autour des listes de numéros associé à des personnes et de lieux. au sommet, un blason :

\cartefigure[0.8\textwidth]
{figures/4 Cartouche Entête.jpg}
{fig:CartoucheCarteBlanc}
{St Maxime de Beaufort
	Le Plan Géographique
	des biens fond que possède le Rd Jean Antoine Blanc tant
	en son chef que par indivis avec Dame Charlotte et Demoiselle
	Jacqueline Blanc ses nièces de même que ceux que
	possède Dame Françoise veuve de Me Michel Blanc
	Notaire en la paroisse de St Maxime
	de Beaufort avec les confins des dits biens
	le tout levé par le géomètre Estienne
	Jaquet sur la mappe originale dressé
	à l'occasion de la mesure générale de Savoye
	[???] 1729 et
	1730}

Nous avons donc un plan dressé par un géomètre, une mention de la Mappe Sarde et plusieurs propriétaires d'une même famille : un prêtre, la veuve d'un notaire et ses deux filles.

Il s'agit clairement d'un objet destiné à être montré et à afficher la richesse et la notabilité de la famille (peut-être en raison de deux jeunes femmes à marier).

Le bourg de Saint-Maxime se situe au centre de la petite vallée du Beaufortain qui était exclusivement rural jusqu'à l'arrivée des sports d'hivers et des barrages hydraulique. Adossée au massif du Mont-Blanc, et entourée des vallées plus large et plus fréquentée de la Tarentaise, de l'Arly. Les sommets qui l'entoure culmine à moins de 3000m, le Doron, gros torrent, s'en échappe vers 800m d'altitude, en laissant les pentes douces avec d'immenses étendues de prairies et de forêts. 

Entre l'isolement et les hivers rigoureux et enneigés, il n'est pas facile d'imaginer la vie dans cette vallée autrement que difficile et pauvre. 
Et pourtant nous avons cette carte ostentatoire.

Pour assoir son emprise directe sur ses sujets et pour mieux contrôler la levée de la taille,  Victor-Amédée II, duc de Savoie (et roi de Sardaigne) met en œuvre entre 1720 et 1730, une "mensuration générale" pour établir un cadastre couvrant toutes les paroisses du duché: c'est la Mappe Sarde à laquelle faisait référence la carte de la famille Blanc.   

Hélène VIALLET a écrit en 1993 un livre passionnant sur l'histoire du Beaufortain, où l'on peut apprendre que la paroisse de Saint-Maxime était une des plus peuplées et une des plus riche du duché de Savoie jusqu'au XVIIIe siècle et cela dés avant l'introduction de la fabrication du gruyère suisse dans la vallée.

Mais ce n'est qu'au cours des années 2000 que le travail de numérisation de la Mappe Sarde par Dominique BARBERO est devenu accessible pour l'historien et  plus de 400 communes de Savoie et de Haute-Savoie ont été ainsi traitées.

La commune de Beaufort fait partie des communes dont la Mappe Sarde a été ainsi numérisée. En exploitant ces données avec les outils dont nous disposons aujourd'hui, nous sommes en mesure d'apporter des éclairages approfondis aux travaux d'Hélène VIALLET sur la paroisse de Saint-Maxime du XVIIIe siècle. En passant, nous essayerons de replacer la carte privée de la famille Blanc dans son contexte.


Nous avons donc pu obtenir de lui une copie de son travail sous la forme de 53 documents pdf contenant le tracé des parcelles tiré de la Mappe originale conservée aux Archives Départementales de Chambéry ainsi qu'un fichier Excel avec les informations des 19791 parcelles relevées par les arpenteurs du Roi de Piemont-Sardaigne.


La stratégie capitaliste de notables investissant dans la propriété des alpages d'altitude a commencé  en Gruyère avec l'apparition du fromage du même nom, par les bourgeois de Fribourg.

L'introduction du Gruyère dans le Beaufortain qui a donné le Beaufort a suivi une logique semblable et si cette évolution est bien documentée pour la famille Blanc dans les travaux d'Hélène Viallet, il nous a paru intéressant d'avoir une vision plus précise des processeurs des alpages : les registres de la Mappe Sarde ne donnent pas l'information sur l'altitude des parcelles. Un travail de géoréférencement a été mené par Dominique Barbéro mais le Cd qu'il a produit sur un grand nombre de communes de Savoie ne donne que des impressions pdf des cartes numérisées. Les Archives de Haute-Savoie en prolongement sans doute des travaux de Dominique Barbéro, a géolocalisé les parcelles pour une partie des communes de Haute-Savoie et permet une consultation en ligne sur leur site, mais ce travail n'a pas été fait par les ADS.
Il m'est donc apparu intéressant de pouvoir calquer la Mappe Sarde sur les cartes IGN et de géoréférencer les parcelles, afin d'identifier d'une manière ou d'une autre les parcelles situées à plus de 1600m d'altitude (nous reviendrons sur ce choix que nous avons fait de cette hauteur). Une fois les parcelles localisées, en extraire celles qui sont en pré ou en pâturage avec le nom des propriétaires devenait faisable en apportant une vraie valeur ajoutée à la lecture des données initiales.

Le roi de Piemont-Sardaigne et duc de Savoie, Charles-Emmanuel III de Savoie décide de permettre le rachat des droits féodaux par les communautés d’habitants par deux édits : l’édit du 20 janvier 1762, puis celui du 19 décembre 1771. Les communautés du Mandement de Beaufort comportant les quatre paroisses : Saint-Maxime de Beaufort, Hauteluce, Villard sur doron et Queige (pour une partie seulement) ont devancé la promulgation de l'édit de 1771 puisqu'elles ont ensemble racheté dés 1770 l'ensemble des droits féodaux les concernant. Le financement de ce rachat a été possible par la vente aux enchères d'une partie des communaux. Dans les documents entourant la procédure des enchères, les parcelles mises en vente sont décrites avec leur numéro du Cadastre Sarde de 1730. Il a donc été possible là aussi de connaître quelles étaient celles concernant des alpages d'altitude.  


	
	
	
\chapter{Contexte}
	
	Avant de se focaliser sur la photographie que la Mappe Sarde offre de la propriété des alpages à un moment précis du XVIIIe siècle (1720-1730), il est utile d'essayer de comprendre le contexte historique et géographique de la vallée du Beaufortain dans une perpective plus globale de la Savoie et des Alpes du Nord, de l'évolution de la propriété  et  l'évolution de l'agriculture et du pastoralisme depuis le moyen-age.
	
\section{Géographie alpine}

a)	Savoie

La Savoie  se situe entre la France, la Suisse et l'Italie, aujourd'hui les grands cols et les grands sommets en tracent la frontière, mais pendant longtemps cette province a constitué une zone à cheval entre le royaume de France et les états italiens.

Les Romains ont créé deux voies romaines. L'une passe par le col du grand Saint-Bernard reliant Genève à Turin, dont le tracé sera repris par la la via Francigena au Moyen Âge pour les pèlerins de Canterbury à Rome.

L'autre, la voie Domitienne, passe par le col de Montgenèvre reliant Milan à Vienne en passant par Turin et Briançon, c'est le passage utilisé par Jules césar pour envahir la Gaule.

Le col du Mont-Cenis (2081m) permet lui aussi de traverser les Alpes de Chambéry à Turin en parcourant la longue vallée de la Maurienne, pour redescendre ensuite le val de Suze . C'est par ce passage que sont passés Charlemagne, Charles le Chauve, l'empereur Henri IV (sur sa route vers Canossa).

La Tarentaise voisine permet aussi un passage par le col du Petit Saint-Bernard et qui aboutit dans le Val d'Aoste. Mais partant de Moutier, il restera plus court pour aller à Turin, de grimper au col de l'Iseran au-dessus de Val-d'Isère pour rejoindre la Haute-Maurienne et le col de Mont-Cenis. Venant du Beaufortain, les chargements de Beaufort et de bétail de la famille Blanc utilisaient cette route via le Cormet de Roselent ou celui d'Arêche.

La vallée de l'Arve qui aboutit à Chamonix en venant de Genève, dispose aussi d'un passage vers le Valais en Suisse par le col de la Forclaz (1527m) mais il faut encore remonter le col du Grand-Saint-Bernard pour se rendre en Italie.

Il est intéressant de noter que la situation exceptionnelle de Turin au débouché de toutes les routes qui traversent l'ouest du massif alpin, a dû être déterminante pour en faire la capitale des États de Savoie en 1563.

Pendant le Haut Moyen Âge, le col de Mont-Cenis fut la route la plus empruntée pour relier les marchands italiens vers les marchés de Bourgogne.
Mais l'ouverture  après 1250 les cols du Simplon et du Saint-Gothard, donnant accès au Rhin va monopoliser les flux commerciaux entre l'Italie (Milan) et l'Europe du Nord. \footcite{PRacine2003} 

xxxxxxxxxxxxxx petite carte !

b)	Particularités du Beaufortain

Le Beaufortain est une vallée en cul-de-sac, elle ne se situe sur aucune voie de communication principale,  son seul accès à  partir de la ville de Conflans (Albertville n'a été créé au XIX siècle) doit emprunter les gorges du Doron et un projet de route est resté en discussion pendant le XVIIIe pour n'aboutir qu'à la fin du XIXe (xxxxxxxxxxx).

Elle est ceinturée de montagnes douces et avec seulement quelques sommets à plus de 2600m. Quatre cols ou cormets permettent d'en sortir : le col des Saisies (1600m) , le col du Joly (2000m, le Cormet de Roselent (2000m) et le cormet d'Arêches (2100m).

Il faut relativiser la notion de cul-de-sac: en effet quand on se déplace à pieds ou à dos de mulet, le passage d'un col ne pose aucun problème en été. Au même titre que les grands cols des Alpes étaient très fréquentés par les voyageurs, les marchands et les armées, les "petits cols" du Beaufortain ont toujours permis une communication facile avec les vallées voisines. C'est seulement en hiver que la vallée devait être complètement fermée sur elle même, vivant alors isolée et en autarcie.

\section{Éléments de contexte historique}



\begin{itemize}
	
	
	\item 1416 : le comte de Savoie, Amédée VIII (\dag 1451), devient le premier duc de Savoie.
	
	\item Le 3 avril 1559 (second traité du Cateau-Cambrésis), le duc Emmanuel-Philibert de Savoie recouvre le duché de Savoie perdu par son père et demeuré sous administration française depuis 1536. Les territoires du duc de Savoie s’agrandissent en direction du Piémont.
	
	\item En 1563, transfert de la capitale du duché de Savoie à Turin.
	
	\item En 1690, l’État savoyard comprend le duché de Savoie, la principauté de Piémont, le duché d’Aoste et le comté de Nice.
	
	\item 1690--1696 puis 1703--1713 : occupations françaises de la Savoie.
	
	\item En 1713, le traité d’Utrecht met fin à la guerre de succession d’Espagne. Victor-Amédée II récupère la Savoie et reçoit la couronne de Sicile. À la suite de l’occupation de la Sicile par les Espagnols, il obtient ensuite le titre de roi de Sardaigne.
	\item 1740 : renversement d’alliance lors de la guerre de succession d’Autriche.
	
	\item 1742--1749 : occupation espagnole.
	
	\item 1792 -- 1815 : la Savoie devient département français du Mont-Blanc
	
	\item rattachement de la Savoie et de l'arrondissement de Nice
	
	
	Au delà de cette petite liste de jalons historiques du duché de Savoie, deux événements ont eu un impact particulier sur la vallée de Beaufort et qui concernent l'objet de notre étude : 
	
	Si l'instabilité et les rivalités locales ont rendu nécessaire la construction d'un grand château-fort par les Seigneurs de Beaufort. Mais 1271, trop endetté, ils vendent les titres et les droits de la seigneurie ainsi que le château lui-même à Béatrice de Faucigny (fille de Pierre II de Savoie et épouse du Dauphin). cet événement est important car depuis cette date et jusqu'en 1681, le mandement de Beaufort restera  rattaché directement aux princes qu'ils soient Dauphinois, Genevoix ou Savoyard.
	
	\item Le Dauphin Humbert II sans descendance vend le Dauphiné au Roi de France en 1349
	\item Le traité de Paris de 1355 met fin aux conflits sans fin entre les Comtes de Savoie , les Comtes de Genève et le Dauphiné. 
	La Savoie obtient une unité territoriale stable et relativement paisible intérieurement. 
	En dehors des occupations militaires françaises et espagnoles de l'époque moderne, la vallée de Beaufort n'est plus une zone frontière à défendre ou à traverser. Le beau château-fort est abandonné et il ne reste aujourd'hui que des vestiges et deux des tours de la forteresse.
	Le pouvoir seigneurial est assuré par des représentants "fonctionnaires" du Duc de Savoie.
	
\section*{Administration communale}


\begin{quote}
	« À partir du milieu du XIII\textsuperscript{e} siècle, les principaux seigneurs de Savoie parviennent à transformer la collection de seigneuries éparses qui constituait jusque-là leur domaine en ce que les historiens nomment une principauté. Il s’agit du comté de Savoie, surtout, mais aussi du comté de Genève et de la seigneurie du Faucigny, contrôlée jusqu’en 1355 par les dauphins de Viennois. Les seigneuries de moindre importance sont vassalisées, y compris celles de l’évêque de Maurienne et de l’archevêque de Tarentaise, ou bien rachetées. Des officiers salariés, les châtelains, sont mis à la tête de circonscriptions fixes, les châtellenies. Les vallées de montagnes qui échappent à la tutelle des évêques et des monastères dépendent ainsi d’un châtelain comtal installé dans un modeste bourg castral : Samoëns, Châtillon puis Cluses, Charousse, Thônes, Faverges, Ugine, Le Châtelard-en-Bauges, Salins, Hermillon puis La Chambre, Beaufort, Entremont-le-Vieux. Des organes centraux de gouvernement sont mis en place : conseil, chancellerie, auditeurs des comptes puis Chambre des comptes. L’administration s’étoffe et développe des procédures écrites de plus en plus sophistiquées. Des enquêtes sont lancées pour connaître les revenus du prince sur son domaine, tandis que les officiers locaux sont soumis à des procédures annuelles de contrôle. Conservés presque intacts jusqu’à nos jours, ces extentes, reconnaissances, et comptes de châtellenie »
	\footcite[59]{Mouthon2016}
\end{quote}


	\begin{quote}
	« Par un édit du 19 septembre 1561, Emmanuel-Philibert permet dorénavant à ces paroisses rurales de s’administrer au moyen d’un conseil élu par l’assemblée générale des habitants et présidé par un ou plusieurs syndics qui en sont les agents d’exécution. Ils ne font pas forcément partie de l’élite sociale, on le voit par exemple dans les registres de pauvreté des Échelles, plusieurs d’entre eux ne savent pas signer et se contentent d’apposer une marque pour authentifier les documents de la gabelle. »
	\footcite[139]{Gachet2011GabelleSel1561}
\end{quote} 


\begin{quote}
	« Dans le Beaufortain, vallée riche en alpages et en forêts s’il en est, les franchises accordées en 1349 à la communauté de Beaufort par le Dauphin alors seigneur du Faucigny, sont confirmées par le comte de Savoie Amédée VI en 1368. Sur dix-sept articles, le contenu se contente de rappeler que le châtelain comtal ne peut concéder les forêts et les pâturages à des personnes étrangères à la vallée. »
	\footcite[65]{Mouthon2016}
	
	\footnotesize
	Note 4 : Auguste Dufour, François Rabut, « Chartes municipales des pays soumis à la maison de Savoie en deçà des Alpes », \textit{Mémoires et documents publiés par la Société Savoisienne d’Histoire et d’Archéologie}, t.~XXIII, 1885, p.~293--311 et p.~325--334.
\end{quote}


\begin{quote}
	« La plus profonde modification des structures administratives au XVI\textsuperscript{e} siècle est liée à la décadence des châtellenies ; nous en avons vu une conséquence dans les difficiles coopérations des châtelains avec les commissaires de la gabelle. Les châtellenies sont à l’origine des circonscriptions administratives, judiciaires et militaires. Mais depuis le XV\textsuperscript{e} siècle, le rôle militaire des châtelains ne cesse de décliner, sauf aux frontières où se maintiennent encore quelques châteaux féodaux. De plus en plus, la défense et la garde des châteaux sont confiées à des fonctionnaires militaires, appelés capitaines ou commandants. S’ajoute à ce dispositif une milice nationale et permanente, créée par Emmanuel-Philibert dans les années 1560. Suite à de nombreuses plaintes pour abus de pouvoir, les prérogatives des châtelains en matière de justice deviennent elles aussi très restreintes. Dans plusieurs châtellenies, le commissaire Ducayre a dû recourir aux forces d’intervention des archers d’un capitaine de justice pour recouvrer des dénombrements que les châtelains refusaient de délivrer. »
	\footcite[138]{Gachet2011GabelleSel1561}
\end{quote}


	\begin{quote}
	« C'est par l'intermédiaire de la commune que le roi communique avec les habitants.
	
	A l'origine, association spontanée de personnes se reconnaissant des intérêts communs, la commune se transforme peu à peu en organisation locale.
	
	L'intervention des membres de la communauté dénommés communiers, s'exerce par l'assemblée des chefs de famille qui sont appelés à désigner un syndic et un conseil, composé de procureurs mandatés pour représenter les communiers dans toutes les affaires les concernant.
	
	Les rois de Piémont-Sardaigne ont donné une existence légale à la commune par l'édit péréquation en 1738.
	
	L'assemblée générale des habitants, a perdu tout pouvoir dans la seconde moitié du XVIII\textsuperscript{e} siècle. Elle n'est plus que réunie, lors d'un fait extraordinaire. C'est ainsi que les communiers s'assemblent pour décider, si oui ou non, il faut s'affranchir des droits féodaux en 1771. En théorie tous les chefs de famille peuvent prendre part à cette réunion, en fait les deux tiers seulement sont présents. Il s'agit des plus influents. Les femmes n'y participent pas, bien que quelques-unes soient responsables d'un foyer. La permission écrite de l'intendant se révèle nécessaire pour convoquer l'assemblée générale.
	
	Les communiers sont représentés par le syndic et le conseil composé de six membres.
	
	Le syndic est recruté chaque année dans le conseil, à l'ancienneté. Rien ne le distingue des autres conseillers, ni son mandat, ni ses fonctions. 
	[...] 
	Au cours des séances, le conseil exprime le plus souvent quelque vœu, "priant très humblement le seigneur intendant" de bien vouloir l'exaucer.
	
	Un délégué transmet en général la requête à l'intendance.
	
	Le procès-verbal de la réunion est dressé par le secrétaire agissant alors en sa qualité de notaire.
	
	C'est lui qui détient tous les pouvoirs de la communauté.
	
	Il exerce de nombreuses fonctions : il dresse les registres divers, établit le rôle de l'impôt du sel, garde dans sa maison, la mappe, le cadastre, les livres de mutation et de délibération. Il tient également la correspondance administrative aussi le choisit-on toujours parmi les notaires, beaucoup plus nombreux qu'aujourd'hui. Il joue dans la communauté le rôle que joue dans la province l'intendant dont il est l'agent direct et responsable. Nommé par cet intendant, sur la proposition du conseil de la communauté, il prête serment, ainsi que le syndic et les conseillers entre les mains du châtelain. Aidé par le secrétaire insinuateur son père (*) jusqu'en 1782, il voit ses pouvoirs constamment renouvelés. De ce fait, on note dans la période étudiée, deux secrétaires seulement. Il s'agit de Joseph Blanc (1750 -- août 1754) et de Michel Blanc son fils (1754--1791) (**).
	
	Ce notaire collégié et secrétaire de la communauté de Saint-Maxime comme il s'intitule, ne laisse que fort peu de pouvoirs au châtelain. 
	
	(*)Le secrétaire insinuateur est chargé d'enregistrer les actes. Pour exercer cette charge à perpétuité et à titre héréditaire il faut disposer de biens équivalent à mille écus d'or. Joseph Blanc achète cette fonction en 1754. Il l'exerce jusqu'en 1782, date de sa mort, il est alors remplacé par son fils Michel.
	
	(**)A partir de 1782, Michel Blanc cumule les fonctions de secrétaire-insinuateur et de secrétaire-notaire. Pour l'aider, en juillet 1784, son cousin, le notaire Jean-Baptiste Blanc, est nommé « exécutant » à l'office de secrétairerie. Cette fonction est révocable. (A.C. BB 3 p. 194).»		
	\footcite[p.~22]{BordeVibertGuigue1971VieRuraleBeaufort}
\end{quote}


\section*{Cadre fiscal}
	


	\begin{quote}
		« La prospérité de 1561 n'avait pas disparu puisqu'en 1574, la communauté put s'affranchir globalement de la taillabilité personnelle envers Mgr de Savoie-Nemours. 543 personnes et leurs familles se libérèrent de ce lien au prix de 16 000 florins et 100 écus d'or, somme payée avec une facilité éloquente par les syndics. »
		\footcite[32]{HViallet1993}
	\end{quote}
	 
	\item Jusqu’en 1559 en Savoie, l’impôt est levé sans régularité, en fonction des besoins et pour une durée limitée. La base de calcul de l’impôt est le feu. L’évaluation globale du nombre de feux et des exemptions accordées aux pauvres permet une répartition par communauté, unité administrative de base.

	\item La gabelle est introduite en Savoie par l’édit du 3 novembre 1560. Il s’ensuit le dénombrement général de la gabelle du sel (1561--1576) en Savoie.	\item Un édit de juillet 1564 crée un droit de subside déterminé selon les ressources. La terre devient alors la base de la nouvelle contribution, qui prend le nom de taille. Les livres fonciers décrivent les confins des parcelles, leur superficie, leur degré de bonté ainsi que le chiffre de la taille. Ils s’accompagnent de livres de mutations, appelés livres de remesses ou de vires. Ces documents sont assortis de cottets, listes ou rôles où figurent les noms des propriétaires et le montant des contributions.\footcite[272]{HVialletSel}
	
	\item L’édit du 24 février 1565 transfère aux communautés la levée de l’impôt : les syndics doivent répartir celui-ci entre les contribuables. Cette réforme marque le début de la promotion de l’institution de la communauté de paroissiens au détriment du mandement ou de la châtellenie.\footcite[272]{HVialletSel}
	
	\item En 1681, Charles-Emmanuel, duc de Savoie, vend et inféode à Joseph-François de Wicardel, marquis de Fleury, la terre, le fief, la juridiction et le mandement de Beaufort --- soit l’ensemble du Beaufortain (Beaufort, Hauteluce, Villard-de-Beaufort et Queige) --- en échange de son palais situé place Royale à Turin (ADC, 208 E-dépôt 2, copie de 1681 d’un acte du 26 décembre 1662). La valeur de cette vente servira ensuite de référence pour évaluer le prix du rachat des droits féodaux.
	
	\item En 1723, Victor-Amédée II promulgue un nouveau code de lois : les \textit{Royales Constitutions}. Il s’agit de l’instauration d’un régime despotique régissant la religion, la magistrature, la législation civile et criminelle, les finances (Chambre des comptes), les droits féodaux ainsi que la gestion du territoire (mines, travaux publics, forêts, etc.). Ces réformes soumettent davantage la Savoie à l’autorité piémontaise et entraînent la suppression de certains particularismes locaux ou féodaux. La réforme la plus importante demeure celle du cadastre.
	
	\item Par lettres patentes du 9 avril 1728, Victor-Amédée II engage les travaux du cadastre général de la Savoie.
	

	
	\item 1762 : première tentative d’affranchissement des droits seigneuriaux (édit du 20 janvier).
	
	\item 1771 : édit du 19 décembre, début de l’abolition des droits seigneuriaux (affranchissement général).
	
\end{itemize}

\section*{Cadre cléricale}

\begin{itemize}
	
	\item Le 26 août 1556, les Beaufortains obtiennent le droit d’élire leur curé et lui versent 300 florins, montant ramené à 200 florins en 1633. Cette coutume perdure jusqu’en 1792 (référence à compléter : \textit{Une belle vallée de Savoie}, J. Garin).
	
	
	
\end{itemize}

\section{la structure de la propriété rurale}


\begin{quote}
	« Entre 1570 et 1625 environ, se produit une transformation des structures foncières. L'indivision disparaît, ou reprend des proportions raisonnables, après le grand affranchissement de la taillabilité personnelle, en 1574. »
	\footcite[92]{HViallet1993}
\end{quote}

\begin{quote}
	« Cette forme de censive, à laquelle nous avons déjà fait allusion, apparaît en Savoie dans la seconde moitié du XIII\textsuperscript{e} siècle, à la faveur de la renaissance du droit romain. Pierre Duparc, dans l’important article qu’il a consacré aux « tenures en hébergement et en abergement », a bien fait le départ entre ces dernières, qui sont, comme nous le savons, des tenures de peuplement concédées à des hôtes chargés de les défricher, et l’albergement emphytéotique. L’un et l’autre contrats ne sont pas sans rapport, comme en témoigne l’identité du vocabulaire : dans les deux cas on dit albergamentum ou abergamentum. Mais en Faucigny, ils ne sont pas contemporains : le second apparaît dans la documentation quand le premier s’éclipse. Les clauses propres à l’albergement emphytéotique montrent qu’il est « le doublet d’un ancien contrat romain qui renaît de ses cendres, un peu modifié, l’emphytéose ». De fait, les documents faucignerans emploient indifféremment les termes « d’albergement », « albergement emphytéotique », ou encore « emphytéose perpétuelle » pour qualifier un contrat par lequel un bailleur (albergateur) concède le domaine utile d’un bien immeuble à un ou plusieurs preneur(s) (albergataire(s)) ; en échange, ceux-ci acquittent le plus souvent un droit d’entrée appelé introge et un cens qualifié de servis, lequel peut être accompagné de prestations secondaires, notamment de tailles et d’aucièges. En Faucigny, l’albergement emphytéotique s’est imposé comme forme unique de tenure au tournant des XIII\textsuperscript{e} et XIV\textsuperscript{e} siècles. Dans les reconnaissances de la fin du XIII\textsuperscript{e}, les tenures anciennement concédées sont encore qualifiées, selon un vocabulaire qu’on retrouve dans toute la Savoie, de « fiefs », et les tenanciers, de « feudataires ». Mais toutes les nouvelles concessions qui se font à ce moment-là sont des albergements. »\footcite[308]{Carrier2001VieMontagnardeFaucigny}
\end{quote}	

\begin{quote}
	« Il arrive même que les pâturages qui font l’objet d’un albergement aient été précédemment tenus en censive. L’alpe de Sales était pâturée par les hommes de divers hameaux de Sixt contre versement d’un auciège, auquel s’était ajouté, à une date inconnue, un cens annuel de 8 sous genevois. Toutefois, l’alpage leur était contesté par les hommes du mandement de Charousse. En 1418, l’abbé alberge à perpétuité l’alpe de Sales, avec ses dépendances des Lanches et du Clouz, à ses hommes de Sixt, contre l’introge considérable de 400 florins, outre le cens et l’auciège qu’ils payaient ci-devant. Le texte de l’albergement est sans ambiguïté : ce que les hommes de Sixt ont acheté si cher, c’est l’assurance de ne jamais se voir préférer ceux de Charousse ; et si les religieux ont renoncé à la liberté d’aliéner l’alpage à qui bon leur semblerait, c’est que le monastère menaçait ruine et qu’ils manquaient de l’argent nécessaire aux réparations les plus urgentes. En l’occurrence, le paiement de l’auciège et d’un cens ne suffisait donc pas à assurer aux hommes de Sixt la quasi-propriété de l’alpage de Sales. C’est qu’au XV\textsuperscript{e} siècle, le droit romain s’est si bien acclimaté en Savoie que l’albergement emphytéotique est devenu la forme normale de tenure. Des droits d’usages même séculaires ne paraissent plus très assurés, s’ils ne sont pas tenus en emphytéose. Ainsi de ceux des consorts de l’alpage de Gers : en 1449, le rénovateur des extentes de Samoëns leur fait remarquer qu’ils ne bénéficient d’un albergement que pour les bois de cette montagne, non pour les pâturages. Il leur faut obtenir du duc la confirmation de leurs droits de pâture, contre un introge de 80 florins. La modicité de la somme, eu égard à la taille de la montagne, s’explique par le fait que l’ancienneté des droits d’usages en question »\footcite[312]{Carrier2001VieMontagnardeFaucigny}
\end{quote}	

	\begin{quote}
	« En sorte que vers 1530, la structure de l’économie du haut Faucigny n’est guère différente de celle qu’on observait vers 1330 : consommer leur orge, leur avoine et leurs raves, manger leur fromage et un peu de viande, dégager quelque argent pour acheter les céréales et le vin qui leur manque, tel est, pour la plupart des paysans du Faucigny, le but de leur travail d’un bout de l’année à l’autre. Il n’y a pas eu de conversion d’une part importante de la population à l’élevage commercial. L’essentiel du commerce de bétail est apparemment fait de petites transactions au jour le jour sur des marchés d’intérêt local. Les foires d’automne donnent lieu à quelques échanges plus importants, où les paysans, en vendant quelques bêtes qu’ils ont engraissées durant l’estive ainsi qu’une partie de leur production fromagère, peuvent acheter avant l’hiver ce que leurs pauvres parcelles n’ont pas su produire et payer leurs termes, dont l’échéance, comme celle des droits seigneuriaux, est généralement à la Saint-Michel. Rares sont les acteurs économiques qui pratiquent un élevage spéculatif : ce sont les plus gros possesseurs de bétail ou, ce qui revient au même, les plus gros possesseurs de prés de fauche, les seigneurs en premier lieu. »\footcite[372]{Carrier2001VieMontagnardeFaucigny}
\end{quote}


\begin{quote}
	« Et d’abord des activités économiques. Les notaires sont généralement de gros possesseurs de terres, et nous les voyons arrondir fréquemment leur patrimoine. L’albergement des tenures échues en est un moyen. En 1484, par exemple, Jean Hoste, notaire, alberge en fief taillable, avec noble Jean Pugin et Guillaume Cochard, les biens échus de Jean Quibliet. L’affermage est aussi un moyen assez fréquemment employé par les notaires ou leurs proches pour acquérir l’usage de terres nouvelles. Nos sources éclairent fort peu ces pratiques, car les princes n’ont pas pratiqué l’affermage à la fin du Moyen Âge. Mais elles sont peut-être plus répendues qu’il n’y paraît au premier abord. En tout cas, les fermiers qui apparaissent dans notre documentation sont toujours des nantis, et fréquemment des notaires, comme les Jay, comme les L’Etelley, que nous évoquerons ci-devant. Il faut aussi mentionner, bien évidemment, les achats purs et simples.
	Pourquoi accumuler ces terres ? Pour en faire des prés de fauche, à n’en pas douter : en 1442, Jacques des Rémondins, notaire, reçoit à albergement une pièce de pré de Pierre Daniel contre l’introge considérable de 200 florins. Le milieu notarial, nous l’avons vu, pratique assez largement le commerce de bétail et de viande. Surtout, on s’en souvient également, il est très largement représenté à la tête des consortages. C’est un autre point à développer : l’omniprésence de la propriété collective ne gêne en rien les gros possesseurs de bétail s’ils contrôlent la gestion des pâturages indivis, ceux des hameaux comme ceux des alpages, c’est probablement que les plus gros troupeaux qui les parcourent leur appartiennent. Leur argent, placé principalement dans l’élevage spéculatif, l'est aussi dans les moulins et autres artifices. »\footcite[488]{Carrier2001VieMontagnardeFaucigny}
\end{quote}


\begin{quote}
	« La grande montagne particulière
	
	C'est le système existant dans le Beaufortain et l'objet même de cette étude. La grande montagne appartient à un particulier, ou à plusieurs membres d'une famille en indivis : elle n'est pas toujours exploitée par le propriétaire lui-même. Le plus souvent, elle est louée pour trois, six ou neuf ans, moyennant une censé annuelle en argent et en denrées à un paysan, le « montagnard », qui travaille avec sa famille et quelques domestiques.
	
	Comme le montagnard est loin de disposer lui-même d'un troupeau suffisant, il prend en location les vaches de son propriétaire et celles des petits éleveurs du Beaufortain, et même de l'extérieur. A l'entrée de l'automne, il restitue les vaches louées et met une partie de son propre troupeau en hiverne chez des paysans de la vallée, qui ont récolté du foin au cours de l'été. Pendant l'hiver, le montagnard redevient un personnage ordinaire avec un rythme de vie et de travail semblable à celui des petits exploitants : consommation du fourrage dans les écuries-granges, production de lait en petite quantité, destiné à l'élevage des veaux et à la fabrication de petits fromages.
	
	L'alpage peut également appartenir à la communauté, mais être exploité de la même façon par un montagnard. La seule différence réside dans l'emploi des redevances du bail : elles sont affectées à l'entretien d'un maître d'école (montagne des Arolles), ou à l'assistance aux pauvres (montagnes de Bellachat et de la Petite Berge). Mais l'acquisition de ces alpages ne fut réalisée qu'à partir du XVIIIe siècle, grâce à des legs privés.
	
	Aucun alpage ne fonctionne selon le système du fruit commun à l'époque étudiée. Le Beaufortain se caractérise donc par une exploitation individualiste, tout en possédant par ailleurs d'immenses communaux : 47 \% du territoire de la communauté selon le cadastre sarde. Cela contraste fortement avec la situation de la Tarentaise, vallée limitrophe. On peut donc se demander à juste titre pour quelles raisons la solution communautaire n'a pas été adoptée. »
	\footcite[48]{HViallet1993}
\end{quote}


\begin{quote}
	« Depuis le Moyen-Age, les communaux rétrécissent comme une peau de chagrin : les syndics ont pour diverses raisons, vendu ou acensé certaines parcelles : mais la cause la plus importante, et la plus grave, est un sournois et irrépressible mouvement d'usurpation des aboutissants.
	
	Dès le XVIe siècle, en raison du poids croissant de la fiscalité, les syndics demandent la permission d'aliéner une partie des communaux pour pouvoir s'acquitter de la taille et des autres charges. Ainsi en 1552 : «pour ce que à présent ledit mandement de Beaufort dèz environ quinze ans en ça, a esté et est fort affollé tant pour raison des tailles et esgances qui ont esté en si grand nombre audit mandement tellement que ce pauvre menu peuple est si fort affollé qu'il ne sçauroit trouver moyen de vivre, veu que le pays est stérille et montaigne, fort habundant en peuple ; désireroyent lesdites leurs communes estre levées en certaines parties es lieus moins dommageables pour le général et icelles estre redduictes en propres... vendre, accenser et admodier à certains particuliers... chose que tombera au grand proffict et utilité de Vostre Excellence et Seigneurie, tant pour raison du direct domaine d'icelles choses... que pour les lodz, ventes et autres obventions que vous surviendront »(A.C.B., DD 1, Requête des communiers, 1552).
	
	A la fin du XVIIe siècle, plusieurs parcelles en montagne furent vendues ou échangées pour permettre l'agrandissement et l'embellissement de l'église de Saint-Maxime. Ainsi le 1er août 1666, Me Claude-Antoine Granier cède «un membre de bouticque proche l'église ; en contr'eschange les sindics lui donnent une montagne de communauté sise aux fins de la Gittaz appelée ez Bancs »(.D.S., Registre 6 E 985, 1666-67). Le même jour ont lieu plusieurs autres ventes, après approbation de l'assemblée des communiers, pour un montant total de 1 200 florins. »
	\footcite[66--67]{HViallet1993}
\end{quote}

\section{les habitants de Saint-Maxime de Beaufort}

c)	Savoie et Alpes du Nord

Les historiens et géographes ont longtemps considérer le pastoralisme comme "inhérent à l'économie alpine". récemment "[...] des historiens médiévistes ont entrepris de réexaminer l’évidence et qu’ils l’ont mise à mal. Ils ont montré que l’élevage n’est pas la vocation naturelle des montagnes, et qu’il n’a acquis la prédominance dans certaines vallées qu’ensuite d’une évolution complexe. Henri Falque-Vert pour le haut Dauphiné, Nicolas Morard pour la Gruyère, Fritz Glauser pour la Suisse centrale et le Tessin et Pierre Dubuis pour le Valais occidental" \footcite{NCarrierEstivage}

« Les historiens helvétiques ont insisté sur les changements que cette évolution a entraînés dans l’occupation du sol. Le cas de la Gruyère est emblématique. Aux XIVe-XVe siècles, la prédominance des ovins y est progressivement remplacée par celle des bovins. En outre, dans le fond des vallées, les prairies gagnent sur les emblavures. Du coup, les prés de fauche qu’on réservait autrefois sur l’alpe sont pâturés. Dorénavant « le plan » et « les monts » sont nettement séparés. Les bêtes ne restent plus, comme il arrivait au XIII\textsuperscript{e} siècle, dans des étables d’altitude pour y consommer le foin récolté sur place. C’est donc au XIV\textsuperscript{e} siècle que l’estivage se met en place tel qu’il existera jusqu’au commencement du XXe , avec l’aller-retour estival entre des fonds de vallée réservés à la culture et à la fauche, et des espaces de pâture se trouvant au-delà de la limite supérieur
de la forêt » \footcite{NCarrierEstivage}

"Un point est assuré : ce ne sont pas les « moines défricheurs » qui sont à l’origine des alpages de la Savoie du nord. D’abord parce qu’ils n’étaient pas assez nombreux ; d’autre part, parce que les textes même des donations dont les religieux ont bénéficié aux Xi\textsuperscript{e}-XII\textsuperscript{e} siècles de la part des grands seigneurs laïcs laissent entendre que les hautes vallées où ils se sont installés n’étaient ni vierges d’hommes, ni inexploitées" \footcite{NCarrierEstivage}

"On voit donc qu’à la fin du XVe siècle, les acteurs de l’activité pastorale dans les massifs considérés sont les mêmes qu’à la fin du XIII\textsuperscript{e}. Si l’on ajoute que de multiples indications laissent à penser que l’élevage reste pour le commun des paysans une activité complémentaire à la céréaliculture, que le rapport entre cheptel ovin et bovin ne paraît pas s’être modifié, qu’on continue comme devant de faucher sur l’alpe" \footcite{NCarrierEstivage}


	\begin{quote}
	« Un document exceptionnel nous donne une photographie de la population 
	beaufortaine en 1561, il s’agit de la Consigne du sel évoquée plus haut. 
	Étudions de plus près le cheptel : le fait le plus frappant est la supériorité 
	numérique des ovins (3\,688 têtes) sur les bovins (2\,890 têtes). Habitués que 
	nous sommes à nos « montagnes à vaches », ce chiffre nous dérange. Les historiens 
	ont pourtant montré que pendant le Moyen Âge, l’arc alpin était abondamment garni 
	de moutons, et que cela ne se limitait pas aux Alpes du Sud. Les troupeaux ovins 
	ont en général moins de vingt têtes ; il s’agit d’une autoconsommation familiale, 
	d’un élevage traditionnel, ne demandant pas de technicité ni de main-d’œuvre 
	nombreuse.
	
	Le troupeau bovin est le plus important économiquement, même si les ovins sont 
	plus nombreux : 2\,890 bovins, 3\,688 ovins. Si l’on divise le nombre de bovins 
	par celui des feux (596), on obtient le chiffre de 4,88 bovins par feu. Si l’on 
	ne prend en compte que les feux possédant du bétail, ce chiffre s’élève à 6,51 
	bovins par feu. Dans la réalité, cette moyenne n’a que peu de valeur car la 
	répartition du bétail est loin d’être uniforme : la majeure partie des paysans 
	(54\,\%) possède de une à dix bêtes ; un deuxième palier regroupe les propriétaires 
	de onze à quinze bovins (13,7\,\% des feux). L’étage supérieur de la paysannerie 
	aisée comprend les 52 feux (13,2\,\%) qui ont à l’étable de seize à trente-cinq 
	bovins. Enfin un dernier groupe, très restreint puisqu’il ne compte que huit feux, 
	possède de 35 à 50 bovins, chiffre considérable. Le schéma démontre clairement 
	que le prétendu égalitarisme montagnard n’existe pas à Saint-Maxime : huit familles 
	sur 596 possèdent à elles seules 12\,\% des bovins. »
	\footcite[90-92]{HViallet1993}
\end{quote}

Nous verrons la part prédominante des communaux en tant que propriétaire des alpages de Beaufort: ils représentent 47\% de la surface de la commune dont 4.600ha soit 70\% en alpage (mais pas forcément des pâturages !!!) xxx vérifier Hviallet p70...xxx

d)	Particularités du Beaufortain

Pour mémoire, le beurre et le fromage ont été inventés par l’homme pour pouvoir conserver dans le temps les calories et la protéine du lait hautement périssable. Cette ressource était d’autant plus nécessaire en zone montagneuse que les cultures de céréales sont incertaines et d’un rendement très médiocre.

Il faut noter l'incroyable imagination contractuelle qui entoure la gestion des propriétés de la terre et du bétail : Imaginons un petit propriétaire d’un terrain de 5ha avec un chalet à une altitude moyenne (1300m) il a plusieurs solutions : il peut avoir 2 ou 3 vaches qu’il fait paître sur une partie de sa prairie, gardant l’autre partie en pré pour faucher 2 fois à la belle saison et faire des réserves de foins pour ses vaches pour passer l’hiver. Avec ses vaches il fera lui-même du fromage (maigre genre petite tomme, ou gras de type vacherin) et du beurre qu’il pourra vendre au marché.

Il avait aussi des ovins et des caprins à hauteur des surfaces en taillis, broussailles et sous-bois qu’il possède ou bien dont il a accès (Communaux).

Il peut aussi garder tout son terrain en pré à faucher, dans ce cas il pourra augmenter le nombre de ses vaches à 5 ou 6. L’été pendant la saison des alpages il pourra laisser sa famille au chalet pour le fauchage, et il mènera ses vaches dans les alpages « communaux » d’altitude. Mais dans ce cas il faudra soit redescendre deux fois par jour son lait à son chalet où il a son matériel pour transformer son lait, soit monter en altitude son matériel à la condition qu'il possède (ou ai la possibilité d'utiliser) un chalet d'alpage.

Une autre possibilité serait pour lui de ne posséder que ses 2 ou 3 vaches et trouver un accord avec un propriétaire plus important à qui il confiera ses vaches l’été et en échange il hébergera l’hiver, deux ou trois vaches de celui-ci en plus de ses vaches qu’il possède.

S’il est vraiment pauvre, il pourra même ne posséder aucune vache en propre, récolter son foin et ne faire qu’héberger l’hiver des vaches d’un propriétaire qui a besoin d’étables où répartir son troupeau l’hiver.

En effet (et jusqu'à la deuxième partie du XXe siècle) il n’existe pas d’étable pouvant héberger plus de 10 ou 15 vaches, alors que les grands troupeaux dépassent les 30 ou 40 vaches. (il faudrait des granges à foin très importantes et du matériel pour amener le foin des prés de fauches à ces granges, puis pour transporter le fumier de l’étable afin d’amender les prés et prairies).

Il serait intéressant de savoir s’il peut faucher ou obtenir une portion de foin récolter sur les prés communaux.

Il doit y avoir des règles d’usages des communaux que ce soit en prés, en alpages (prairies) ou en bois pour les petits propriétaires.

Peut-être existe -t – il des chalets d’altitudes communs, où chacun des utilisateurs des alpages peut y garder son matériel (ou utiliser un matériel commun, ou faire le beurre et le fromage en commun et se le repartir ensuite (à « fruit commun » ?

Un plus gros propriétaire pourra avoir plusieurs chalets à différents étages, il fera progressivement sa « remue » partant du chalet le plus bas au printemps et déplaçant son troupeau vers le haut au fur et à mesure de la fonte des neiges (on trouve encore des éleveurs au XXIe siècle qui « remuent » entre cinq ou six chalets.

\section {citations non utilisées}

	\begin{quote}
		« De fait, en Dauphiné comme en Savoie, les mas, qui désignent à la fois des portions de territoire et les collectifs de tenanciers qui les exploitent, comprennent non seulement des parcelles de « terre » (des champs), des prés, voire des maisons et des jardins mais également de larges portions de l’espace inculte, essentiellement des bois et des pâturages. Lors de l’enquête de 1265 les habitants de Césane, en Briançonnais d’Outre-Monts (piémontais depuis 1713), rappellent que, depuis longtemps, leur territoire paroissial est partagé en mas et que ceux-ci s’étendent depuis l’eau qui coule au milieu du territoire jusqu’au sommet des montagnes. »
		\footcite[50]{Mouthon2019MontagnesMedievales}
	\end{quote}
	


\begin{quote}
	« Deux actes d'état de 1681(A.D.S., registre de Me Guyard 6E992, 1680-1681) illustrent bien le système des remues. Ils sont d'autre part très intéressants, car on dispose pour ces deux alpages du calendrier de leurs remues, pratiquement identiques, à l'époque contemporaine.
	
	Hoirie de Me Jean Doix(cte d'état pour Anne Missilier veuve de Me Jean Doix, 5 octobre 1681, registre 6 E 992, fol. 584.): montagne du Cormet. [...]. Cet alpage présente une importante dénivellation : 600 m séparent la station de base (village de Roselend, 1 475 m) et le chalet du Cormet (2 050m).
	
	En 1681, il existe huit bâtiments, dont deux en très mauvais état :
	— à Roselend, une grande maison et un grenier,
	— le bâtiment de la Saulsaz (1 645 m),
	— le bâtiment du Plan de la Lai (1 830 m),
	— le bâtiment du Gollet,
	— le bâtiment de la Garettaz, en ruine,
	— le bâtiment du Cormet (2 027 m),
	— le bâtiment du Lard (2 050 m), en ruine.
	
	[…]
	
	L'étude de l'hoirie de Noël Frison(A.D.S., registre de Me Guyard 6 E 992, fol. 555), décédé en 1662, nous montre la répartition des terres et pâturages d'une famille de montagnards : ses fils Jean et Maxime, qui se partagent les biens de leur père, disposent d'une exploitation dans le bas de la vallée, à Domelin (850 m), mais la maison familiale est aux Curtillets (960 m). L'été venu, ils remuent vers la vallée de la Gittaz, où se trouvent leurs montagnes, après une étape à la Tournaz (1 370 m).
	
	Quelques années plus tard, selon un acte d'état, il existe sur les alpages huit bâtiments dont deux en médiocre état :
	— au village de la Gittaz, deux maisons comprenant cuisine, sertour, freidier, grenier, grange, étable (1 665 m),
	— deux granges au sommet du village,
	— le bâtiment des Exerts (1 840 m),
	— le bâtiment de la Costaz (1 910 m),
	— le bâtiment des Crétets (2 180 m),
	— le bâtiment des Jovets (2 240 m) »
	\footcite[54--55]{HViallet1993}
\end{quote}



	
	
		\begin{quote}
		« Bien qu’ayant très tôt concédé l’usage des pâturages aux paysans, les seigneurs, tant laïcs qu’ecclésiastiques, n’ont donc pas pour autant renoncé à l’exploitation directe en matière d’élevage, à l’exception (notable) des comtes de Genève et de Savoie. Pour les uns, les alpages dans lesquels on les voit intervenir semblent être des éléments de leur réserve, à l’exploitation de laquelle ils ont associé leurs dépendants, sans l’abandonner tout à fait. C’est que les surfaces dont ils disposaient étaient bien trop importantes pour qu’ils puissent les mettre en valeur eux-mêmes, mais qu’ils ne voulaient pourtant pas renoncer à posséder un cheptel, pour leur consommation personnelle ou dans un but commercial. C’est ainsi que lorsqu’il alberge la montagne de Sales en 1418, l’abbé de Sixt se réserve le droit d’y inalper chaque année cinquante génisses et génissons avec les bêtes des albergataires. Concédant la montagne des Chéserys en 1444, le prieur de Chamonix garde la possibilité d’y inalper autant de bêtes laitières ou non laitières que les autres consorts. Dans un tel cas de figure, le seigneur perçoit l’auciège — au sens de redevance — à la fois parce qu’il est le détenteur du domaine éminent et parce qu’il fournit les bâtiments ; mais il possède également une part du domaine utile, au même titre que les autres consorts et étant soumis aux mêmes contraintes, comme le précise clairement l’albergement des Chéserys. En ce cas, on peut supposer, même si les textes ne l’indiquent pas, que s’il a inalpé des laitières, il a sa part de l’auciège, entendu cette fois au sens de production fromagère totale. Mais nous voyons aussi des nobles posséder des droits d’usage dans les alpages communaux dont ils ne détiennent pas le domaine éminent. C’est même, semble-t-il, le cas le plus fréquent. En 1538, par exemple, noble Georges, fils d’André de la Frasse, reconnaît tenir du duc de Savoie, en fief noble : « La moitié d’une cabane située à Monjover. Une cabane située aux prés aux paquiers communs. Une autre aux communes d’Harmancette. Le droit que luy et ses hommes ont sur les joïres, paquiers et communes existants depuis le nant d’Albuergant jusqu’a celuy d’Hontares. L’aucheage qu’il a coutume de percevoir sur ses hommes ». \footcite[344--345]{Carrier2001VieMontagnardeFaucigny}
	\end{quote}	
	
	

	
\chapter{Le Beaufortain}

\begin{table}[htbp]
	\centering
	\caption{Les 20 villes et bourgs les plus importants de Savoie (nombre d’habitants)\footcite[37]{Gachet2011GabelleSel1561}}
	\label{tab:villes_savoie_1561_1776}
	\resizebox{\textwidth}{!}{%
		\begin{tabular}{lrr|lrrr}
			\toprule
			\multicolumn{3}{c|}{En 1561} & \multicolumn{4}{c}{En 1776} \\
			\cmidrule(r){1-3} \cmidrule(l){4-7}
			Ville & Popul. 1561 & Popul. 1776 & Ville & Popul. 1561 & Popul. 1776 & Évol. \% \\
			\midrule
			Megève            & 5057   & 2473 & Chambéry         & 5002   & 9755 & +95\% \\
			Chambéry          & 5002   & 9755 & Annecy           & 2871   & 4478 & +56\% \\
			Thorens           & 4457   & 1637 & Thonon           & lacune & 2901 & -- \\
			Beaufort          & 3429   & 2796 & Beaufort         & 3429   & 2796 & -18\% \\
			Samoëns           & 3410   & 2573 & Samoëns          & 3410   & 2573 & -24\% \\
			Bourg-Saint-Maurice & 3164 & 1607 & Saint-Jean-de-M. & 1960   & 2474 & +26\% \\
			Sallanches        & 3113   & 1184 & Rumilly          & 1539   & 2273 & +48\% \\
			Chamonix          & 2962   & 865  & Saint-Pierre-d’Albigny & 1695.8 & 2243 & +32\% \\
			Saint-Gervais     & 2942   & 1722 & Yenne            & lacune & 1952 & -- \\
			Annecy            & 2871   & 4478 & Thônes           & 2447   & 1947 & -20\% \\
			Taninges          & 2748   & 1602 & Saint-Martin-de-B. & 1759 & 1863 & +6\% \\
			Thônes            & 2447   & 1947 & La Roche         & 1240   & 1860 & +50\% \\
			Faverges          & 2280   & 1803 & Faverges         & 2280   & 1803 & -21\% \\
			La Motte          & 2265   & 1511 & Ugine            & 2235   & 1803 & -19\% \\
			Ugine             & 2235   & 1803 & Moûtiers         & 1216   & 1801 & +48\% \\
			Passy             & 2204   & 1514 & Saint-Gervais    & 2942   & 1722 & -41\% \\
			Meillerie         & 2070   & 1167 & Valloire         & 2007   & 1648 & -18\% \\
			Valloire          & 2007   & 1648 & Grand-Bornand    & 1412   & 1643 & +16\% \\
			Saint-Jean-de-M.  & 1960   & 2474 &                  &        &      &      \\
			\bottomrule
		\end{tabular}
	}
\end{table}

Remarques : le titre du tableau est ambigu car il compare des entités urbaines (Chambéry, Annecy...) avec des paroisses rurales où la population indiquée n'est pas celle du bourg principal mais celle de la paroisse tout entière(Beaufort, Megève...).



\section{écosystème clos ? unité pastorale et géographique particulière}
	-	Pré pâturage\\
	-	Culture\\
	-	Bois\_
	-	Altitude moyenne\\
	-	Grande superficie\\
	-	Une seule voie tout temps\\
	-	4 cols d’altitudes\_
	-	Stabilité population\\
	-	Grande Montagne et forêts\_
	
	Dates de la Saint-Jean à la Saint Martin\_
	Cormet Turin via petit saint-bernard 212klm\_
	Cormet Turin via col de l’iseran mont Cenis 186klm (Rivoli 169klm)\\
	Beaufort Turin via Albertville et maurienne 214klm (Rivoli 197klm)\\
	
	Gabel du sel beaufortain 1561 2890 bovins / 3688 ovins => pas inférieur à ceux du 19e ? et préférence bovine
	
		\begin{quote}
		« Les paroisses qui recensent des bêtes étrangères sont exactement celles qui comptent des bêtes en hiver, avec une exception notoire, celle de Beaufort. C’est la paroisse qui hiverne le plus de bétail, mais il n’y a pas une seule bête étrangère recensée ! Elles ne sont pas souhaitées car les espaces communaux sont surpâturés et usurpés, soit entre communautés soit par des paroisses voisines. Pour y mettre des bêtes il faut normalement être communier et ne pas en prendre plus qu’on ne peut en hiverner. Ce règlement est régulièrement enfreint, comme s’en plaignent les syndics et procureurs en 1552 : « …très humblement exposent les scindiques et procureurs du mandement de Beaufort, vray estre comme en ladite communauté sont plusieurs communes lesquelles se despacquent, gaudissent, dissipent et affolent par gens non point contribuables aux charges dudit mandement… et pour ce que plusieurs sont venus achepter des montagnes près et adjacentz ausdites communautés, lesquels ne sont aucunement habitantz duddit mandement ny moing font aulcun support esdits communiers… » [...] L’absence totale de bêtes étrangères dans le recensement de 1561 semble montrer que le règlement était finalement bien appliqué. »
		\footcite[88]{Gachet2011GabelleSel1561}
	\end{quote}
	
\section{Exception féodale et cléricale ?}
	-	Directement dépendant du duc de Savoie puis du marquis de Fleury\\
	-	Aucun résident grand propriétaire\\
	-	Pas de couvent ni monastère\\
	-	Dépendance (lâches ?) avec Saint Maurice (Valais) puis archevêque de Moutiers\\
	-	Choisi ses curés\\
	-	Pas de contrats d’albergement\\
	-	Superficie des communaux\\
	
	La communauté de Beaufort
	
		\begin{quote}
		« En Savoie, le cas de l’église de Beaufort est, de ce point de vue, 
		l’exception qui confirme la règle. Au XVe siècle, la communauté souhaite 
		depuis longtemps imposer à la tête de l’église Saint-Maxime un desservant 
		originaire de la vallée. Elle s’oppose en cela au patron légal de l’église, 
		c’est-à-dire la très ancienne et prestigieuse abbaye valaisanne de 
		Saint-Maurice-d’Agaune, qui a succédé dans ce rôle à l’archevêque de 
		Tarentaise. En 1514, une transaction aboutit à officialiser la désignation du 
		curé de Beaufort par les syndics, les conseillers et les principaux chefs de 
		famille de la communauté. Le postulant devait en outre obligatoirement être 
		du pays, suivant la formule \textit{ecclesia Bellifortis nescit extraneos} : 
		l’église de Beaufort ne veut rien savoir des étrangers. [...]
		
		Par leur maîtrise de l’écrit, ces prêtres contribuent à renforcer le prestige 
		mais aussi le rôle d’intermédiaire avec l’extérieur de ces familles dirigeantes. 
		Ils monopolisent aussi les fonctions de chapelains. Ils desservent les chapelles 
		de villages, notamment celles fondées par leurs familles, et surtout disent les 
		messes payées par les testateurs de leurs villages et des familles alliées. »
		\footcite[148--149]{CarrierMouthon2010PaysansAlpes}
	\end{quote}
	
		\begin{quote}
		« La noblesse savoyarde est essentiellement urbaine. Malgré l’homogénéité de la répartition de la population, elle n’investit pas pour autant les paroisses montagnardes, excepté peut-être dans le Faucigny. Elle se rapproche donc plus de la noblesse italienne issue du patriciat marchand des villes que de la noblesse française d’origine féodale qui a une prédilection rurale. »
		\footcite[140]{Gachet2011GabelleSel1561}
	\end{quote}
	
\section{Le fromage de Beaufort}
	
		À un moment, entre la fin du Moyen Âge ou vers le XVIe siècle, de gros propriétaires ont appris à faire des fromages différents, à pâte cuite, plus gras et riche que la tomme (qui est écrémée afin de faire aussi du beurre), qui peuvent se conserver plus longtemps (2 ou 3 ans !). Ce fromage exigeait des volumes de lait important (technique de fabrication et/ou optimisation de l’usage du bois pour chauffer les chaudrons ?).
	
	Ce fromage nouveau pouvait se vendre nettement plus cher que les fromages traditionnels et que le beurre. L’économie pastorale s’est transformée en conséquence, et il devenait sans doute plus intéressant à notre petit propriétaire de confier tout son troupeau à un gros propriétaire et de profiter marginalement de la plus-value apportée par ce nouveau fromage.
	
	Ce nouveau fromage exclusivement fabriqué à partir de lait de vache a orienté l’agriculture à se spécialiser dans l’élevage bovin au détriment des ovins et caprins ( ??? stat avant et après des troupeaux).
	
	Le mythe des moines défricheur (Mouthon communs)
	
		\begin{quote}
		« C’est clair en tout cas, pour la Gruyère : il s’agit d’une montagne 
		principalement ovine à la mi-XIV\textsuperscript{e} siècle, et l’on y pratique 
		la polyculture vivrière ; trois siècles plus tard, c’est une montagne à vaches 
		toute tournée vers la fabrication et la commercialisation de grosses meules de 
		fromage à pâte dure. Nicolas Morard place les prémices de ce changement au 
		XV\textsuperscript{e} siècle, et le mouvement des enclosures lui en paraît le 
		signe le plus indubitable. De fait, la volonté de récolter plus de foin peut 
		aisément être mise en parallèle avec un essor de l’élevage des bovins, qui en 
		sont, comme chacun sait, de gros consommateurs. Mais il ne faut pas oublier la 
		modestie de cette première vague d’enclosures. Quant au commerce du fromage, 
		c’est une date symbolique pour la région que la première mention, en 1449, de 
		l’achat « d’ung fromagiierre de Gruery » par l’hôpital Notre-Dame-des-Bourgeois 
		de Fribourg. Toutefois, on a en Gruyère comme ailleurs des traces d’un petit 
		commerce du fromage et des autres produits de l’activité pastorale depuis le 
		XIII\textsuperscript{e} siècle, et ni le volume ni l’ampleur des transactions 
		ne paraissent changer d’échelle au XV\textsuperscript{e}. Le pas décisif vers 
		la spécialisation de la Gruyère vers l’élevage et la commercialisation du fromage 
		se situe lorsque la Ville et République de Fribourg annexe le comté entre 1536 
		et 1555. Elle reconnaît alors comme conforme « au droit divin, humain et naturel » 
		que les cultivateurs puissent jouir de leurs prés en toute propriété. Ils 
		obtiennent donc un droit de clore total et sans limite, avec cette justification 
		que par les enclosures « n’est fait tort à ceux qui ont moins de possessions veu 
		que du mesme peu qu’ilz ont, ilz tireront la jouissance entière ». Dans les 
		décennies qui suivent, des bourgeois de Fribourg achètent en Gruyère de grandes 
		quantités d’alpages, et en réalité ce sont eux qui vont tourner l’ensemble de 
		la région vers l’élevage commercial. »
		\footcite[287]{CarrierMouthon2010PaysansAlpes}
	\end{quote}
	
	\begin{quote}
		« Nous savons maintenant que l’histoire de l’élevage n’est pas immobile. Hélène Viallet a montré par exemple que la fabrication des premiers gruyères de Beaufort n’est pas antérieure aux années 1630-1640, qu’elle a été permise par l’appel à des fruitiers suisses venus de la région de Gruyère, et qu’elle est à l’origine d’une mutation de l’organisation pastorale, aboutissant à la généralisation du système de la « grande montagne particulière. »\footcite[317--319]{Carrier2001VieMontagnardeFaucigny}
	\end{quote}	
	\begin{quote}
		« La plus ancienne mention du gruyère, appelée aussi « grevire »(38)se trouve dans une quittance du 6 juin 1638 : «vingt florins pour reste de la vente de fromage en forme de gruiere... » (A.D.S., E 203, Registre du notaire Bochage, 1638). La fabrication reste d'abord limitée à quelques alpages. Partout ailleurs, on fait le fromage traditionnel, la « robbe d'esté » dont nous ne connaissons pas les caractéristiques ; il s'agissait certainement d'un fromage à pâte pressée cuite. En 1641, coexistent «fromage gruiere» et «fromage de pays», mais à partir de 1650 on ne trouve plus cette dernière catégorie dans les redevances en nature. »
		\footcite[96]{HViallet1993}
	\end{quote}
	\begin{quote}
		« Cinq types d’intervenants ont part aux charges comme aux profits résultant de l’exploitation d’un alpage, sachant que concrètement, des personnes peuvent appartenir à deux ou plusieurs de ces catégories : 1) le ou les détenteurs du domaine éminent, c’est-à-dire le seigneur ; 2) le ou les détenteurs du domaine utile, c’est-à-dire les consorts ; 3) le ou les possesseurs des chalets, qui peuvent ou non faire partie d’une des deux catégories précédentes, voire des deux à la fois ; 4) le personnel qui accompagne les bêtes à l’alpage pendant le temps de l’estivage et qui fait le fruit ; 5) des propriétaires de bétail qui, sans avoir part au domaine utile de la montagne, auront confié quelques bêtes à l’un des consorts en sorte qu’il les inalpe pendant l’été. Cette pratique, appelée « collecta animalium », donnait lieu à des contrats comme celui par lequel Pierre Addam, curé de Notre-Dame du Lac, prend en location cinq vaches de Jean Vernuz du 24 juin au 8 septembre 1460. Il lui en coûte quatre florins et quatre gros, en échange de quoi il aura la totalité du lait de ces vaches pendant la période en question. Ces dates, qui encadrent les deux mois d’été, indiquent que le curé compte faire inalper ces vaches, dans le but de réaliser un bénéfice sur le fruit réalisé avec leur lait. C’est le seul exemple que le minutier Solliard ait conservé de ces accords, qui paraissent pourtant avoir été assez fréquents, comme on le verra. Outre le loyer qu’il perçoit, qu’on peut imaginer aussi à part de fruit, le propriétaire des bêtes y gagne d’en être débarrassé pendant les mois d’été et de n’avoir pas à les faire paître sur ses prairies du bas, qu’il peut réserver à la fauche. Mais il n’a nulle autre part aux charges et aux revenus de la montagne. »\footcite[330]{Carrier2001VieMontagnardeFaucigny}
	\end{quote}	
	
	\begin{quote}
		« le maître fruitier se trouve au sommet de la hiérarchie de la domesticité montagnarde. Fier de son art, il en fait un secret
		jalousement préservé. Ce sont les fruitiers du pays de Gruyère, engagés par les	notables [...] Dès 1646, plusieurs baux prévoient que l'ascensateur de la montagne paiera le salaire du fruitier. « Grueren », « grue » deviennent synonymes de fromager. Travailleur qualifié, le fruitier régente le personnel de l'alpage. Le montagnard lui-même doit toujours le consulter et tenir compte de ses avis. C'est en fait lui qui a la responsabilité la plus lourde, puisque, de la vente du fromage, et donc de sa qualité dépendent le paiement du loyer et un éventuel bénéfice. C'est pourquoi il exige des gages très élevés : 52 £ 10 s. en 1773 pour le fruitier fribourgeois de la communauté de Sainte-Hélène-sur-Isère, qui « s'oblige de venir faire la gruyère de l'année prochaine dans la montagne appartenante à la communauté, d'y faire le beurre, y faire faire le sérat, et de diriger la montagne en homme de bien et d'honneur, et de la même manière que tous suisses ont coutume de faire ... et que la gruyère soit bonne et recevable, et salée à propos »(G. PÉROUSE, La commune en Savoie. Introduction à l'inventaire sommaire des archives communales de l'arrondissement d'Albertville, 1911., p. LXXIII.). Me Joseph Blanc, qu'on ne peut taxer de prodigalité, donne 70 £ à son fruitier Gaspard Balmaz pour la campagne de la même année»
		\footcite[148]{HViallet1993}
	\end{quote}
	
	\begin{quote}
		« Pendant le reste de l'année, on ne fabrique pas de gruyère. En effet, personne ne dispose d'un troupeau assez important pour obtenir trois cents litres de lait par jour, le minimum pour pouvoir faire un gruyère. Les montagnards et les paysans qui possèdent un nombre suffisant de vaches font des tommes grasses appelées parfois « boudanes »<35',ou maigres si l'on a ôté la crème du lait pour faire du beurre. Les tommes de Savoie sont à pâte pressée non cuite, affinée au moins un mois en cave froide et humide. [...] La fabrication du fromage n'emploie que 10 \% du lait, il reste 90 \% de petit-lait dont on va tirer le beurre de fleurettes, puis le sérac. La matière grasse restant en petite quantité est récupérée par le chauffage du petit-lait à une température de 75-80 °C, après addition d'aisy <40'.Elle monte en surface sous forme de petites fleurs que l'on recueille avec une cuillère plate. Après refroidissement pendant vingt-quatre heures, elles sont barattées. Le beurre de fleurettes est bien entendu moins riche que celui élaboré à partir de crème
		fraîche. [...] (le) vacherin pèse deux à trois kg. C'est
		un fromage gras et coulant, « lorsqu'il est un peu attendu, il se sert à la cuillère 	et s'étend comme de la crème ». Très apprécié des amateurs de bonne chère, il figure en bonne place sur les tables de la noblesse et de la bourgeoisie»
		\footcite[187-191]{HViallet1993}
	\end{quote}
	
	
\chapter{Constitution du corpus cartographique et cadastral}

Décision du Duc de Savoie : 1er cadastre européen\\
10 ans / geomètre\\

\section{Les sources mobilisées}

L'étude repose sur trois ensembles documentaires complémentaires.

Le premier est constitué de la Mappe sarde de Beaufort conservée aux Archives départementales de la Savoie sous la cote 208 Edépôt 25-61. La documentation comprend les deux toiles cadastrales dressées au début des années 1730 ainsi que les registres associés : livres des numéros suivis, tabelle-minute, tabelle générale et registres de calculs. Ces documents permettent de connaître pour chaque parcelle son propriétaire, sa nature, sa superficie, son degré de bonté et diverses informations fiscales.

Elle se compose de plusieurs documents graphiques (toiles peintes) dont la version en couleur est disponible sur le site internet des ADS \footnote{Cote C 2155 Mappe 361 et 361bis : Copie de la mappe originale (plan cadastral) de la commune de Beaufort. Signée Cocelli, directeur du bureau de la péréquation pour le duché de Savoie, 1732, plan au lavis, papier collé sur toile , dimension 712x328 et 700x260 cm \url{https://recherche-archives.savoie.fr/?id=recherche_guidee_cadastre_web}}

Les archives numériques proposent 4 photographies pour une vue générale (2 par toiles) ainsi que 45 photos de vue détaillée.

Les ADS possèdent aussi les registres associés à la Mappe : 

- Cote C 2154 : Mappe originale (plan cadastral) de la commune de Beaufort, dans lequel sont figurées et numérotées les 17691 parcelles qui composent le territoire de la commune. Signé Gastaldy et Tirola, géomètres

-Cote C 2156 à C 2160 5 volumes : Livres des numéros-suivis du géomètre (tableau indicatif) de la commune de Beaufort, contenant, par ordre de numéros-suivis, le numéro de chacune des parcelles figurées à la mappe, le nom du propriétaire, le lieu-dit et la nature de culture. 

- Cote C 2161 : État des erreurs reconnues sur la mappe et dans les livres du Cadastre de la commune de Beaufort, dressé par le géomètre qui a opéré les corrections. - (En italien.)

- C 2164 : Tabelle-minute (cadastre primitif) de la commune de Beaufort, contenant les noms des propriétaires par ordre alphabétique, les numéros sous lesquels les parcelles qui leur appartiennent sont inscrites à la mappe, le lieu-dit, la nature de culture, la contenance en mesures de Piémont et en mesures de Savoie. - (Il manque, au commencement du registre, une partie de la lettre A.) - Biens du marquis de Beaufort ; de noble J. -B. de Burgard ; des religieux cordeliers de Sainte-Marie Égyptienne de Chambéry ; du marquis de Cordon ; de la cure de Beaufort ; des chapelles des Pénitents de Beaufort, d’Arêches, de Notre-Dame des Châteaux, de Notre-Dame de la Pitié, à Dommelin, du Bersend, de Beaubois, de Beaudin, des Charnets, de la Frasse, de la Portella, de la Pierre, des Houtards, des Curtillets, du Mont, du Prat, de Roselen, de la Gitte, des Villes-Dessus, des Villes-Dessous, de la Dray ; du marquis de Saint-Innocent ; de la Communauté de Saint-Maxime de Beaufort ; du village d’Arêche ; du bourg de Beaufort ; des villages de Beaubois, du Bois, du Bersend, du Chapet, du Contafolliat, des Curtillets, des Charnets, des Houtards, du Planay, de Pesay, de la Pierre, du Prat, de Roselen, des Traverses, de la Toile, de la Serra, des Villes-Dessous, de Gravier, du quartier du Pré. - Surface cadastrée de la commune : 50324 J. 208 T

C 2162 : Livre de calculs de la commune de Beaufort, contenant, par ordre de numéros-suivis, les calculs des surfaces de tous les triangles qui composent chacune des parcelles, ainsi que la contenance de chaque parcelle en mesures de Piémont et de Savoie.

- 2e volume, du n° 6380 à 10844. - Le 1er volume manque
C 2163 : Livre de calculs de la commune de Beaufort. - 3e volume, du n° 10852 à 17691. - État d’erreurs rectifiées et de calculs nouveaux au sujet des parcelles atteintes par l’inondation.
C 2165 : Brouillard des déductions de la commune de Beaufort, contenant tous les calculs pour les déductions qui ont eu lieu pour diverses causes, ainsi que ceux de la cote de la taille.
C 2166 à C2172 (7 volumes) : Tabelle générale (cadastre récapitulatif) de la commune de Beaufort, contenant, outre les éléments du cadastre-minute, le degré de bonté, les augmentations et déductions pour revenus, frais de culture, servis, etc., et enfin la cote pour la taille



Les ADS n’ont pas fait le travail de numérisation et géocodage tel qu’à pu le faire les Archives Départementales de Haute-Savoie. Pour les communes du Beaufortain et alentour, les ADS ne disposent que des documents originaux : les informations de la Mappe et des registres associés ne sont ni numérisés ni géolocalisés et encore moins retranscrit par OCR.

Une exception toutefois pour Beaufort, Dominique Barbero en amont du travail fait par la suite par les ADHS, à fait un travail disponible en CD-Rom sur 432 communes de Savoie et de Haute-Savoie dont la commune de Beaufort mais malheureusement pas cette des autres communes du Beaufortain ni des communes limitrophes.

Le second ensemble est donc constitué du travail réalisé par Dominique Barbero. Celui-ci a produit pour la commune de Beaufort un atlas numérique comprenant cinquante-cinq cartes sectorielles de la Mappe ainsi qu'un fichier de données descriptives reprenant les informations cadastrales de l'ensemble des parcelles. Ce travail constitue une ressource exceptionnelle car aucun équivalent n'existe pour les autres communes du Beaufortain.

Le CD-Rom de Beaufort est disponible en consultation aux ADS ainsi qu’à la vente directement auprès de Mr Barbéro.

Le CD-Rom contient : 
-	Une présentation de l’Atlas
-	55 images contenant chacune un secteur de la Mappe avec le contour des parcelles et leur numéro.
-	Un fichier Excel de données attributaires tirées des registres avec un descriptif par parcelle \footnote{le fichier Excel est aussi disponible sur le site des ADS}.
-	Des cartes diverses et données de propriétaires 

Nous ne savons par quels moyens et outils a effectué ce travail . Les cartes fournies en format image (.jpg) sont clairement issues d’un logiciel GIS. 
En revanche, la géolocalisation des parcelles n’est pas fournie dans le CD-Rom.

Notre objectif poursuivi par le présent travail fut  de géolocaliser la Mappe de Beaufort, puis les parcelles disponibles sur les images disponibles du CD Barbéro avec le logiciel QGIS. 

Enfin, les photographies numériques mises en ligne par les Archives départementales de la Savoie ont été utilisées afin de contrôler les informations fournies par l'atlas Barbero et de revenir aux documents originaux lorsque cela était nécessaire.

\section{Limites des données disponibles}

Malgré leur richesse, ces sources présentent plusieurs limites.

Les registres cadastraux ne sont pas numérisés sous une forme directement exploitable pour l'analyse statistique. Les cartes disponibles sont des images raster et non des données géographiques vectorielles. Les informations descriptives et les informations cartographiques sont donc dissociées.

Par ailleurs, toutes les parcelles décrites dans les registres ne sont pas directement localisables sur les cartes fournies par Dominique Barbero. Un travail préalable de reconstitution et d'enrichissement du corpus a donc été nécessaire avant toute analyse historique.

\section{Constitution d'une base de données enrichie}

Afin de rendre possible l'exploitation historique du cadastre, les données initiales ont été complétées par plusieurs informations nouvelles.

Les graphies manifestement variantes d'un même mas ont été harmonisées afin de constituer des ensembles territoriaux cohérents. Un attribut \texttt{Mas2} a ainsi été créé pour faciliter les traitements statistiques et cartographiques.

Les nombreuses dénominations cadastrales de la rubrique « Nature » ont également été regroupées dans une nomenclature simplifiée d'occupation du sol. Cette opération avait pour objectif de conserver la richesse descriptive des registres tout en permettant une lecture synthétique de l'organisation du territoire.

Enfin, chaque parcelle a été associée à une altitude obtenue par croisement avec un modèle numérique de terrain. Des classes altitudinales ont ensuite été définies afin de permettre l'étude de la répartition des propriétés, des usages du sol et des formes d'exploitation selon l'altitude.

\section{Résultat obtenu}

À l'issue de ce travail, chaque parcelle dispose :

\begin{itemize}
	\item des informations cadastrales issues des registres sardes ;
	\item d'une localisation géographique ;
	\item d'une appartenance à un mas ;
	\item d'une catégorie homogène d'occupation du sol ;
	\item d'une altitude et d'une classe altitudinale ;
	\item d'une superficie exprimée en hectares.
\end{itemize}

Cette base de données constitue le corpus utilisé dans l'ensemble des analyses présentées dans les chapitres suivants.

Les traitements techniques ayant permis la constitution de ce corpus sont présentés en détail dans l'annexe méthodologique.

\section{Méthodologie}

\subsection{Analyse de la table du CD-Rom : numérotation des parcelles}

Le CD-Rom Barbéro contient un ficher Excel : Beaufort-FICHIER-ATLAS-OK.xls contenant 17803 parcelles  (avec en plus, une ligne d’en-têtes des colonnes) :

•	Numéro : il s’agit des numéros de parcelles a priori numériques dont le plus grand numéro est 17691. Il devrait a priori avoir 17691 parcelles sur la commune de Beaufort du Cadastre Sarde mais on constate immédiatement que le nombre de lignes ne correspond pas au nombre de parcelles

•	CF : Chef-lieu « O » ou « null » ; signification à investiguer ultérieurement

•	E : Ancien fond ecclésiastique « O » ou « null » ; signification à investiguer ultérieurement

•	I : « O » ou « null » ; signification à investiguer ultérieurement (dans la notice livrée avec le CD-Rom il est fait mention à cet emplacement d’une colonne nommée « N » (pour Noble)

•	Propriétaires : nous avons là des noms et prénoms des propriétaires avec parfois des compléments d’information par exemple « MONOD Jeanne soit Popelin veuve d'Antoine GACHET » ou « TARIN Nicolas et Humbert son frère» et aussi des noms de communautés par exemple « La communauté pour l'usage commun » ou « Chapelle de Roselend »

•	Statut : cette colonne contient du texte avec des valeurs répété : « Bourgeois » « communier » … il s’agit d’une information sur le statut du propriétaire a priori tel qu’il est inscrit sur le registre associé au cadastre

•	IND : parcelle en indivision, elle contient des valeurs « null » et quelque « O »

•	Mas : cette colonne contient du texte avec des valeurs répétées : « Roselend » « Lay » … ce sont des noms de lieux, le mot « Mas » correspond souvent à un hameau ou à un lieu-dit de la commune ou plutôt toute la zone territoriale d’influence de ce hameau

•	Nature : cette colonne contient du texte avec des valeurs récurrentes : « Roc », « Pâturage » ; « Marais et pâturage » … qui sont des descriptions de la nature du sol

•	DB : cette colonne contient des valeurs numériques 0, 1 , 2 ou 3 qui correspond aux quatre « degrés de bonté ». Le degré de bonté est une évaluation de la qualité (ou du rendement) du sol, faite par les agents du duché de Savoie en charge du relevé du cadastre, il a un but fiscal

•	m² : il s’agit sans doute de la surface de la parcelle en m² ; le registre du cadastre contenant de valeur de superficie selon une double mesure : la mesure de Savoie (Journaux, toises, pieds) et mesure de Piémont (Journaux, Tables, Pieds). 1 journal de Savoie équivaut à 0,3685 ha, celui du Piémont lui est légèrement supérieur. Cette conversion a été réalisée sans doute par le créateur du fichier D.BARBERO. (remarque : l’ensemble de la commune actuelle couvre 149,53 klm²  mais la partie transférée à l’époque de l’établissement de la Mappe Sarde est revenue à la commune actuelle de Beaufort sur Doron, la somme des surfaces du fichier D.BARBERO est égale à 139.648.325 m² soit 139,6 klm² ce qui est cohérent, on peut donc valider que la surface sur fichier Excel est en m² d'autant plus qu'à l'époque de l'établissement de la Mappe une partie des parcelles a été rayée car réattribuée à la paroisse de LaChapelle, elles ont été depuis réintégrées à la commune de Beaufort, dont il est normal que la superficie calculée par Barbéro soit inférieure à la superficie de la commune actuelle de Beaufort-sur-Doron) 

•	Avoine : des valeurs null ou numérique s’étalant de 13,5 à 22

•	Seigle : des valeurs null ou numérique s’étalant de 0,5 à 4 (les valeurs 0,5 sont stockées au format texte « 1 / 2 » nous les avons changées en 0,5

•	F. Bœuf : sous bœuf est regroupé le nombre de bovins, la colonne contient des valeurs numériques de 0,5 à 9 puis pour les valeurs supérieures le nombre est suivi de la lettre « L » : « 10 L », « 12 L », « 100 L » …

•	F. Cheval : sous Cheval est regroupé le nombre d’équidés de toute sorte (âne, mule, mulet cheval ) . La colonne contient des valeurs numériques de 0,5 à 9 puis pour les valeurs supérieures le nombre est suivi de la lettre « L » : « 10 L », « 12 L », « 100 L » …

•	Fascines : (fagot de petit bois ?) valeur numérique non investiguée pour l’instant

•	Chevron : valeur numérique sans doute bois d'œuvre pour charpente

•	Latte : valeur numérique sans doute bois d'œuvre pour des planches

•	Pièce 1 : valeur numérique non investiguée pour l’instant

•	Pièce 2 : valeur numérique non investiguée pour l’instant

•	N° Grief : champs pouvant contenir une ou plusieurs valeurs numériques, il s’agit de numéro de référence des contestations faites auprès des autorités relativement à la parcelle cadastrée. En principe il faudrait normaliser ce champ et créer une logique de relation 1-n entre une parcelle et les griefs associés, cependant il n’y a que 22 cas où une parcelle est associée à 2 ou 3 griefs, et l’étude des griefs est très secondaire dans notre étude.

•	Grief : champ texte contenant une typologie selon les types de griefs

•	Problèmes : descriptif sommaire du grief


Degré de Bonté :

Dans l’attribution des degrés de bonté, chaque type d’occupation du sol comporte les trois classes ; la qualité de la parcelle pouvant être bonne (1), moyenne (2) ou mauvaise (3). Par exemple pour la paroisse de Landry en vallée de Tarentaise (Savoie) la « note des estimes en degré avec leurs catégories… » nous informe sur la production qu’un terrain doit fournir pour se voir attribuer tel ou tel degré de bonté. Un pâturage doit produire annuellement deux quintaux de foin pour obtenir le degré de bonté 1, un quintal et demi pour le 2 et un quintal pour le 3. Pour les champs, la mesure est en bichet et il en faut vingt-neuf répartis en froment, seigle et farine pour répondre aux critères du degré de bonté 1, vingt-six pour le 2 et vingt-trois pour le 3, et ainsi de suite pour chaque usage du sol. Pour le bâti, les estimateurs attribuent le degré de bonté d’une terre voisine, de préférence appartenant au même propriétaire. Cette méthode d’évaluation des terrains fixe des niveaux d’exigence précis basés sur le rendement de la parcelle. Ainsi, les parcelles classées en 1 sont de bonnes terres productives, car les critères exigés pour avoir ce degré de bonté sont du même niveau. Le fait que chaque nature de sol comporte l’intégralité des classes (1, 2 et 3) et que les critères d’évaluation répondent à un même niveau d’exigence permet une comparaison entre deux terrains du même degré de bonté. \footcite{Baud2010} 


Mesure de Savoie :
« La base de ce système est le « pied de chambre de 0,033m qui élevé au carré donne le « pied de cadastre de 0.92m². un carré de 8 pieds forme la toise carrée (7.370m²) et 400 toises carrées équivalent à un journal de Savoie dont la valeur exacte est de 29.4838 ares. On compte pour les approximations 3 journaux à l’hectare » 	\footcite{GUICHONNET1955}


Mesure de Piemont :

« Les cadastreurs piémontais se servirent des unités de leur pays et qui dérivaient du pied liprand de 0,513 m. Il fallait 6 pieds pour faire un trabuc (3,082 m) et 2 trabucs pour une perche (6,164 m). Ces mesures linéaires servaient de base aux unités agraires suivantes : le pied de table (rectangle de 1 x12 pieds liprands de côté) ou 3,167 m². La table, ou perche carrée vaut 12 pieds de table (38,009 m) et 100 tables, ou 400 trabucs carrés, donnent un journal de Piémont. La mesure de Piémont est à préférer lorsqu’on utilise l’ancien cadastre. Pour les usagers soucieux de l’exactitude des contenances, elle est évidemment meilleure car elle résulte des opérations sur le terrain ; pour l’historien, elle est commode car la conversion des tables en journaux repose sur la base 100. » Idem


Mas :

«On a la preuve, par les reconnais; féodales, que les mas sont antérieurs de trois ou quatre siècles au cadastre, et que les « trabucants » piémontais n’ont fait qu’enregistrer un état de fait profondément et obscurément enraciné dans la tradition paysanne. Les mas appartiennent à plusieurs propriétaires, nobles ou roturiers, et le régime foncier de 1730 recouvre une réalité agraire très antérieure; ils chevauchent parfois plusieurs communes et ne portent pas toujours le nom des villages ou des lieux-dits qu’ils renferment, mais un toponyme d’étymologie romaine ou germanique. Leur étendue comprend rarement une seule zone de végétation, mais englobe différents types de terroirs et semble bien représenter une unité élémentaire d'exploitation, domaine d’un seul individu. » Idem p296

a) Nettoyage des données

•	Pb colonne « Numero »

En fin de fichier il y a 13 lignes avec des valeurs alphanumériques : 12 parcelles ont une lettre à la suite du numéro de parcelle et 1 parcelle a une lettre sans numéro :

17160 B;;;;BOUCHAGE Jean Baptiste;..\\
17161 A;;;;BUGAND Claude;… \\
17301 A;;;;MARTIN Claude;… \\
17302 B;;;;MARTIN Claude;… \\
17303 C;;;;La communauté pour l'usage commun;Communaux;… \\
17304 D;;;;La communauté pour l'usage Commun ;Communaux ;… \\
17305 E;;;;La communauté pour l'usage commun;Communaux;…\\
17306 F;;;;La communauté pour l'usage commun;Communaux;… \\
17307 G;;;;La communauté pour l'usage commun;… \\
17308 H;;;;La communauté pour l'usage commun;…\\
17309 I;;;;La communauté pour l'usage commun;\\
17310 K;;;;La communauté pour l'usage commun;\\
A;;;;TOCCAT Daniel;Communier;;;;3;1543;;;;;;;;;;;;\\

Nous avons créé une nouvelle colonne : Complément_numéro.

Les lettres des 13 lignes sont déplacées dans cette nouvelle colonne, et la ligne sans numéro se voit attribuer un numéro « 0 ».

Il y a aussi des parcelles dont le numéro est suivi de 1/2 , 1/3 ou 1/4 , ces valeurs ont été remplacées par des décimales (0.5, 0.3, 0.25) ajoutées au numéro de parcelle.

Il reste à comprendre pourquoi le nombre de lignes est significativement différent de la valeur maximum de numéro de parcelle. Mais nous allons d’abord importer ce fichier dans une base de données de travail.

•	Pb colonnes « CF » , « E », « I » , « IND »

Les valeurs « O » dans les colonnes « CF » , « E », « I » , « IND » ont été remplacées par des « Oui » afin de permettre une caractérisation booléenne.

•	Pb colonne « m² »

147 ;;;;VIALLET Jacques fils de Benoit;Communier;;Praz;Chenevier;1;\#VALUE!;21;3;;;;;;;;;;\\
La valeur \#value est remplacé par la valeur « null » \\
•	Pb colonne « Nature »

La colonne Nature est longue et contient beaucoup de double description, par exemple « bois et pâturage » ou « prés et teppes », une option de normalisation aurait été de décomposer en une relation 1-N : une parcelle pouvant être reliée alors à deux natures « bois et « pâturage » ; mais n’étant pas certain de comprendre précisément l’intention du géomètre-cadastreur, nous avons pris l’option de laisser pour l’instant le champ nature tel quel, mais en revanche de créer un concept « Occupation_Sol » pour regrouper les nombreuses « Natures » en ensemble homogène en prenant pour principe lorsqu’une Nature est « multiple » de prendre la première comme la Nature principale ; par exemple « « bois et pâturage » » aura la valeur « Bois » pour « Occupation_Sol »

•	Pb colonne « Mas »

Pour les Mas il y a des mas dont les libellés se ressemblent, par exemple : « Coutafalliat » et » Coutassalliat » (les » s » et les «f » dans la graphie du XVIIIe), mais comme les numéros de parcelles ne sont pas forcément continus entre les deux orthographes, nous avons décidé de laisser les libellés de mas tel quel en attendant un travail futur de représentation cartographique des mas et parcelles.

•	Pb colonnes « Avoine », «Seigle », « Boeuf », « Cheval », « Fascines » , « Chevron » , « Latte» :

Les champs « Avoine », «Seigle », « Boeuf », « Cheval », « Fascines » , « Chevron » , « Latte» décrivent soit une quantité de production d’Avoine de seigle, de fagots de bois (Fascines ?) et c’est logique qu’elle soit associée à une certaine parcelle précise, mais pour les autres il s’agit d’un nombre de bêtes ; cette dernière valeur pose question, car cela devrait:

•	Être associé au propriétaire et correspondre à la taille de son troupeau (cela poserait alors d’autres questions, car pour la gestion des alpages (en été) les troupeaux sont constitués de bêtes appartenant en propre et des bêtes en location) ; cela ne semble pas être le cas, car pour un propriétaire possédant plusieurs parcelles le nombre de bêtes varie selon une logique non claire ; exemple :
Parcelles	Propriétaires	Mas	Nature	DB	Surface	Boeuf\\
7806	BLANC Jean Baptiste	Val de Treicol	Grange	2	67	180\\
7807	BLANC Jean Baptiste	Val de Treicol	Grange	2	181	180\\
7808	BLANC Jean Baptiste	Val de Treicol	Pâturage	2	31130	180\\
7774	BLANC Jean Baptiste	Val de Treicol	Grange	2	206	160\\
7776	BLANC Jean Baptiste	Val de Treicol	Pâturage et rocher	2	86210	160\\
8331	BLANC Jean Baptiste	Couecle	Pâturage et broussailles	2	65791	150\\
8323	BLANC Jean Baptiste	Couecle	Pâturage	2	147294	120\\
8324	BLANC Jean Baptiste	Couecle	Grange	2	79	120\\
7802	BLANC Jean Baptiste	Val de Treicol	Pâturage	3	10135	50\\
7803	BLANC Jean Baptiste	Val de Treicol	Pâturage	3	22085	30\\

•	Ou être associé à une ou plusieurs parcelles : il est évident qu’un troupeau qui « remue » change de parcelle régulièrement.
Nous avons donc conservé ces champs tels quels comme attribut de la parcelle en attendant de comprendre comment cette donnée fonctionne (il faudra sans aucun doute revenir aux registres originaux manuscrits associés au Cadastre Sarde conservé aux Archives Départementales de Savoie)

•	Pb colonnes « Avoine » et « Seigle »

Par exemple dans la ligne :

5299 ;;;;DUITTOZ GENIS;Communier;;Bersend;Champ;3;377;14 1/2; 1/2;;;;;;;;;;\\
Il est inscrit 14 1/2 pour 14,5 et 1/2 pour 0.5

En ouvrant le fichier avec Excel, les valeurs sont bien comprises par le logiciel comme des nombres décimaux, sauf lorsque la valeur est 1/2 qui reste comprise comme du texte

Les 668 cas des colonnes « Avoine » et « Seigle » ont été corrigés.

•	Pb colonnes Bœuf et Cheval

Il existe des « L » qui suivent les valeurs supérieures à 10 ; je ne sais pas à quoi cela correspond (sans doute pas l’unité de mesure Livre. Ces « L » sont supprimés afin d’avoir des colonnes numériques. (Il faudra tout de même vérifier ce qui est inscrit dans le registre.
Il y a aussi un cas où la cellule Bœuf contient une virgule ; elle est remplacée par « Null»

Pour la colonne Cheval il y a aussi 4 cas de « 1/2» remplacés par « 0.5»

•	Pb colonnes No Grief, Grief, Problème

Il existe des lignes dont le champ « Problèmes» n’a pas de numéro de Grief associé sans numéro., nous avons ajouté une numérotation allant de 800 à 952 et la description « Grief sans numéro ».

•	divers :
Nous avons aussi modifié certains titres de colonnes (remplacement des é , ², œ, espace…) et ajouté une colonne » Analyse » (vide) qui servira pour stocker les résultats des analyses des numéros de parcelles.

Désormais nous avons un fichier csv contenant des champs propres et homogènes :

Analyse	Texte\\
Numero	Numerique\\
Complement_numero	Texte\\
CF	booléen (Oui/Non)\\
E	booléen (Oui/Non)\\
I	booléen (Oui/Non)\\
Proprietaires	Texte\\
Statut	Texte\\
IND	booléen (Oui/Non)\\
Mas	Texte\\
Nature	Texte\\
DB	Numerique\\
Surface	Numerique\\
Avoine	Numerique\\
Seigle	Numerique\\
Bœuf	Numerique\\
Cheval	Numerique\\
Fascines	Numerique\\
Chevron	Numerique\\
Latte	Numerique\\
Pièce1	Numerique\\
Pièce2	Numerique\\
No_Grief	Texte\\
Grief	Texte\\
Problemes	Texte\\

\subsection{Constitution d’une table attributaire enrichie}

Un travail de modélisation de la table contenue dans le CD Barbero est disponible en annexe: ce travail permet de distinguer les entités : Mas, Propriétaire et Nature

Au niveau des Mas nous avons par exemple fusionné le mas "Coutassallat" dans le mas "Coutafaillat" ainsi que le mas "Treverse" dans le mas "Traverse"

Au niveau des nombreuses "natures", nous avons créé une entité "Occupation du Sol" afin de regrouper les 440 natures différentes.

Regroupement des natures cadastrales

Les registres de la Mappe sarde décrivent chaque parcelle au moyen d'une rubrique dite « nature ». Cette information est particulièrement riche et reflète la diversité des milieux exploités par les communautés montagnardes du Beaufortain. Outre les catégories classiques telles que les champs, prés ou bois, les estimateurs distinguent de nombreuses situations intermédiaires : pâturage sur roc, roc et pâturage, bois et broussailles, rocher broussailles, glières, rippes, marais, rivages, etc. Au total, plusieurs dizaines de libellés différents apparaissent dans les registres.

Cette richesse descriptive constitue un atout pour comprendre la perception et l'évaluation des terres par l'administration sarde. En revanche, elle rend difficile toute analyse statistique ou cartographique globale. La multiplication des catégories produit des tableaux peu lisibles et empêche de dégager les grandes structures de l'occupation du sol à l'échelle de la communauté.

Afin de permettre les traitements statistiques, les différentes natures ont donc été regroupées en un nombre limité de catégories fonctionnelles. Ce regroupement ne vise pas à remplacer les dénominations originales, qui demeurent conservées dans la base de données, mais à faciliter l'analyse des grandes composantes du territoire.

La méthode retenue repose sur une hiérarchie fondée sur l'usage économique principal de la parcelle. Lorsqu'un libellé comporte le terme « pâturage », la parcelle est classée dans la catégorie « Pâturage », quelle que soit la présence éventuelle d'autres éléments (roc, bois ou broussailles). Ce choix s'explique par l'importance économique de l'activité pastorale dans le Beaufortain et par le fait que de nombreuses parcelles d'alpage associent naturellement herbe, affleurements rocheux et végétation ligneuse. (remarque : nous n'avons pas estimé fiable l'idée que l'estimateur sarde ait marqué une prédominance avec l'ordre des môts dans son descriptif, ainsi "pâturage et bois" et sera pas différentié de "bois et pâturage").

Les parcelles ne comportant pas de pâturage sont ensuite examinées pour identifier les terres destinées aux cultures (champs, jardins et autres terres cultivées), puis les prés de fauche. Viennent ensuite les espaces boisés, puis les broussailles et bruyères. Enfin, les terrains essentiellement minéraux (roc, roche, rocher, ravines, murgers, glières, etc.) sont regroupés dans la catégorie « Roc ». Une catégorie spécifique est également conservée pour le bâti ainsi que pour les surfaces en eau, marais et rivages. Cette logique hiérarchique est biaisé et orientée avec le parti pris de a) l'importance économique majeur des pâturages pour le Beaufortain en 1730, puis b) les champs permettant des cultures vivrières et c) les prés de fauche de fourrage pour l'hiver, nous avons placé d) les bois avant e) les broussailles (ces dernières ne sont pas du tout improductives car elles permettent de nourrir des caprins et même de fournir une alimentation en hiver pour tout le bétail ) enfin nous avons f) les roc g) les marais... h) le bâti

Cette classification présente l'avantage de restituer les grandes fonctions économiques du territoire tout en restant fidèle à la logique des registres sardes. Elle permet notamment de distinguer les espaces agricoles, les ressources fourragères destinées à l'hivernage du bétail, les pâturages d'altitude, les ressources forestières et les terrains à faible potentiel productif. L'analyse croisée de ces catégories avec l'altitude et le degré de bonté permet ainsi la mise en évidence de l'organisation agropastorale de la communauté de Beaufort au XVIIIe siècle.


	
\chapter{Analyse du Territoire de Beaufort}

\section{Généralités}

La commune actuelle de Beaufort qui correspond exactement à la paroisse de Saint-Maxime de Beaufort du 18e siècle (à l'exception d'un territoire de 830 ha objet d'un litige avec la paroisse des Chapelles cf.plus loin).

L'altitude minimale (730m) correspond au fond de la vallée du Doron en direction de Villars-sur-Doron et culmine à 2882m. 

Elle était à l'époque moderne la paroisse la plus peuplée et celle payant le plus d'impôt du duché de Savoie 

(xxxxAjouterSourcexxxx)

Il s'agit de la deuxième plus grande commune de Savoie et de Haute-Savoie après celle de Bourg-Saint-Maurice (les communes actuelles de Val-Cenis et LesBelleville sont, aujourd'hui, administrativement plus grandes , en raison des regroupements de communes récents). La taille moyenne en 2024 des communes françaises était de 1500 ha.

La superficie actuelle est de 14.927 ha (15.276 ha mesurés par QGIS sur la superficie de la couche commune de l'IGN)
dont 10.135ha (\texttt{66\%}) se situent au-dessus de 1600m
et 2.482ha (\texttt{16\%}) sont composés de pentes supérieures à 40°

\cartefigure
{figures/1_ZonesPentues40.jpg}
{fig:ZonePente}
{zones à forte pente (>40°) et zones en altitude (>1600m))}
 
 Une des caractéristiques du Beaufortain réside en son paysage doux avec des sommets moyennement élevés, peu abrupts et avec des prairies qui s’approchent de la plupart des sommets.
remarque : Les 137 parcelles dépourvues d'informations attributaires ont été exclues des traitements statistiques. Elles correspondent à des parcelles initialement portées sur la mappe de Beaufort puis rayées et réattribuées à la paroisse des Chapelles. Ne conservant ni nature de culture, ni propriétaire, ni données fiscales exploitables, elles ne peuvent être intégrées aux analyses quantitatives. 


Rappel : la notion d’occupation du sol est un attribut créé lors de notre présent travail. Il n’existe pas dans les registres de la Mappe Sarde ni dans le travail de Dominique Barbéro. Il regroupe des « Natures » de sol qui sont trop nombreuses pour être facilement compréhensibles.


(\texttt{38\%})









\section{Altitude, occupation du sol et organisation du territoire}
	
	L'analyse croisée de l'altitude, de l'occupation du sol et de la répartition des biens communaux met en évidence une organisation verticale particulièrement marquée du territoire beaufortain. Les différentes formes de mise en valeur ne se répartissent pas de manière homogène mais s'inscrivent dans un système d'étagement qui structure l'ensemble de l'espace communal.
	
\subsection{Une organisation verticale du territoire}
	
	
	
	\begin{table}[htbp]
		\centering
		\includegraphics[
		width=\textwidth,
		trim=0cm 8cm 0cm 0cm,
		clip
		]{tables/1 Surface par altitude et OccSol.pdf}
		\tabcaptionindex{tab:occsol-alt}{Occupation du Sol et Altitude : Surface en ha}
	\end{table}
	
	Nous retrouvons par ce tableau l’image d’un territoire très important  composé principalement de pâturage (\texttt{38\%}), de forêts (\texttt{16\%}) de champs et de prés (\texttt{12\%}) et de rocs (\texttt{24\%}). 
	
	La répartition des occupations du sol selon l'altitude fait apparaître une rupture nette aux environs de 1600 mètres.
	
	Sous cette altitude se concentrent la quasi-totalité des champs et jardins ainsi qu'une très forte proportion des prés. Les cultures céréalières et les espaces de production fourragère se développent principalement entre 800 et 1400 mètres, là où les conditions climatiques demeurent les plus favorables.
	
	À l'inverse, au-dessus de 1600 mètres, les surfaces cultivées deviennent marginales tandis que les pâturages occupent une place prépondérante. Les bois, broussailles et terrains rocheux représentent également une part importante des surfaces recensées. Cette évolution traduit les contraintes imposées par l'altitude mais également l'adaptation des activités humaines aux différents étages de la montagne.
	
	La rupture observée vers 1600 mètres ne correspond toutefois pas à une limite entre un espace occupé et une montagne déserte. Elle marque plutôt le passage entre deux formes complémentaires d'exploitation du territoire : d'une part un espace dominé par les cultures et l'habitat permanent, d'autre part un espace largement consacré aux activités pastorales.
	
\subsection{Une montagne intensément fréquentée}
	
	L'étude du bâti confirme cette interprétation.
	Les parcelles "bâties" sont plus intéressantes à analyser en nombre de parcelles plutôt qu'en superficie.
	
	\begin{table}[htbp]
		\centering
		\includegraphics[
		width=\textwidth,
		trim=0cm 7cm 0cm 0cm,
		clip
		]
		{tables/1_Alt_OccSol_Bati.pdf}
		\tabcaptionindex{tab:occsol-bati}{Analyse des Bâtiments (en nb de parcelles)}
	\end{table}
	
	Les maisons sont majoritairement concentrées dans les étages inférieurs, particulièrement entre 800 et 1400 mètres d'altitude. Cette répartition correspond à l'implantation des principaux villages et hameaux permanents de la paroisse.
	
	Cependant, les hautes altitudes ne sont nullement dépourvues de constructions. Au-dessus de 1600 mètres, les registres recensent encore un nombre important de bâtiments, principalement des granges, des greniers, des maisons d'alpage et des masures. Leur présence témoigne d'une occupation saisonnière intense de la montagne.
	
	Cette répartition reflète l'organisation traditionnelle de l'économie beaufortaine. Les troupeaux sont conduits vers les alpages durant la belle saison et la fabrication du fromage s'effectue directement dans les montagnes. Les bâtiments d'altitude constituent ainsi les infrastructures indispensables au fonctionnement du système pastoral : logement temporaire des hommes, abri des animaux, stockage des produits agricoles et transformation du lait.
	
	L'opposition entre les étages altitudinaux ne doit donc pas être comprise comme une séparation entre espace habité et espace vide, mais comme la coexistence de deux formes d'occupation complémentaires : un habitat permanent concentré dans les basses et moyennes altitudes et un habitat saisonnier étroitement associé à l'exploitation des alpages.
	
\subsection{Les communaux : une propriété essentiellement pastorale}
	
	L'analyse des biens communaux révèle leur forte concentration dans les étages supérieurs de la montagne.
	
	\begin{table}[htbp]
		\centering
		\includegraphics[
		width=\textwidth,
		trim=0cm 7cm 0cm 0cm,
		clip
		]
		{tables/1_Superf_Communaux.pdf}
		\tabcaptionindex{tab:SurfCommun}{Surface des Communaux en ha}
	\end{table}
	
	
	Les communaux sont composés principalement de pâturages, de bois, de broussailles et de terrains rocheux. À l'inverse, ils comprennent très peu de champs, de prés ou de bâtiments.
	
	Cette répartition ne doit pas être interprétée comme le maintien collectif de terres marginales abandonnées par les propriétaires privés. Les espaces d'altitude constituent au contraire le support d'activités essentielles à l'économie locale. Les grands pâturages, les parcours de montagne et les ressources forestières représentent des biens dont l'exploitation nécessite souvent une gestion collective et des règles d'usage communes.
	
	La concentration des communaux dans les hautes altitudes semble ainsi traduire la permanence d'un modèle ancien de gestion communautaire adapté aux contraintes de l'économie pastorale. Loin d'être périphériques, ces espaces occupent une place centrale dans le fonctionnement économique de la communauté.
	
\subsection{La question du degré de bonté}
	
	\begin{table}[htbp]
		\centering
		\includegraphics[
		width=\textwidth,
		trim=0cm 10cm 0cm 0cm,
		clip
		]
		{tables/1_OccSol_DB.pdf}
		\tabcaptionindex{tab:occsol-db}{Occupation du Sol par Degré de Bonté (en \% de Surface)}
	\end{table}
	
	
	Les biens communaux présentent une proportion particulièrement élevée de parcelles classées dans le meilleur degré de bonté.
	
	Cette observation mérite cependant d'être nuancée. Le degré de bonté utilisé dans la fiscalité sarde ne constitue pas une mesure strictement agronomique de la fertilité des sols. Il s'agit avant tout d'un outil d'évaluation fiscale destiné à apprécier la valeur économique des biens afin de répartir l'impôt.
	
	Dans une région où l'élevage et l'exploitation des alpages jouent un rôle déterminant, les estimateurs ont pu attribuer des degrés élevés à des terrains dont l'intérêt économique était important, même lorsque leur aptitude à la culture restait limitée. Un alpage productif pouvait ainsi recevoir une évaluation favorable sans pour autant présenter les qualités agronomiques d'une terre labourable.
	
	Il faut également tenir compte de la logique de péréquation fiscale mise en œuvre à l'échelle du duché. Afin de maintenir une répartition cohérente de l'impôt entre des communautés aux ressources très différentes, les agents chargés des estimations ont vraisemblablement accordé une attention particulière à la valeur économique des ressources pastorales qui constituaient l'une des principales richesses du Beaufortain.
	
	Les degrés de bonté élevés observés pour les communaux doivent donc être interprétés avec prudence. Ils témoignent probablement davantage de la valeur économique reconnue aux alpages, pâturages et espaces forestiers que d'une qualité agronomique exceptionnelle des sols.
	
\subsection{Conclusion}
	
	L'étude croisée de l'altitude, de l'occupation du sol, du bâti et du statut foncier met en évidence une forte spécialisation verticale du territoire beaufortain.
	
	Les basses et moyennes altitudes concentrent l'habitat permanent, les cultures et les principaux espaces de production fourragère. Les hautes montagnes sont quant à elles dominées par les pâturages, les espaces forestiers et les terrains de parcours indispensables à l'économie pastorale.
	
	Cette organisation ne traduit pas une opposition entre espaces riches et espaces pauvres ni entre territoires occupés et délaissés. Elle révèle au contraire la complémentarité de différents étages de montagne dont chacun remplit une fonction spécifique. Les communaux occupent une place centrale dans ce système en assurant la gestion collective d'une part importante des ressources pastorales et forestières qui fondent la prospérité économique de la communauté.
	
	


\section{Questionnement sur les mas et les patronymes}

	\begin{quote}
	« Gabriel Pérouse écrivait qu’en Savoie, les communautés paysannes sont avant tout des associations de propriétaires, tenant en indivision des biens concédés par un seigneur. C’est absolument vrai en Faucigny, à condition de préciser que les communuitates de hameaux, les plus nombreuses, sont probablement assez souvent, à l’origine, des communautés lignagères. De nombreux hameaux faucignerans ont donc une origine familiale. Ils sont nés et se sont développés parce que des cognats continuaient à voisiner après avoir quitté la cohabitation, et gardaient sur de nombreuses générations une partie de leurs biens en communs, ainsi devenus communia de hameaux. La cohésion familiale que nous avons constatée aux XIV\textsuperscript{e} et XV\textsuperscript{e} siècles vient donc à la suite d’une longue tradition. Elle s’explique notamment par la dureté des conditions de mise en valeur du sol aux XII\textsuperscript{e}-XIII\textsuperscript{e} siècles.»\footcite[248]{Carrier2001VieMontagnardeFaucigny}
\end{quote}

	\begin{quote}
	« En 1377, il y a bientôt un siècle que les descendants des Rey tiennent en indivision le pâturage jadis accensé par la Grande Dauphine, et ils ne manifestent nullement l’intention de se séparer. D’où la tentation de penser que les communaux de hameaux, qui sont si nombreux aux XIV\textsuperscript{e} et XV\textsuperscript{e} siècles, sont, pour une bonne partie d’entre eux, d’anciennes propriétés familiales restées indéfiniment indivises. D’autres indices viennent confirmer cette idée. À Chamonix, on s’en souvient, les mas et les « prises » restent des unités de perception encore vivaces au XV\textsuperscript{e} siècle, alors qu’on ne les voit plus guère mentionnés dans les autres mandements. Seule une petite partie d’entre eux sont devenus des hameaux. Mais que ce soit ou non le cas, tous ceux qui subsistent au XV\textsuperscript{e} siècle ont leurs communaux. Or, ce sont certainement, avons-nous dit, d’anciennes exploitations familiales. Il semble donc bien que les communaux de prises et de mas correspondent à la part des terres de ces exploitations qui n’a jamais été partagée. Dernier point enfin, si certaines familles tiennent leur nom d’un village, nombreux sont les villages qui ont pris le nom d’un lignage : Les Loyers, Les Echenaz, Les Meuniers, Les Julliards, Les Sauvages, Les Curtets, Les Bollus, pour ne prendre que quelques exemples, sont autant de hameaux qui correspondent à des anthroponymes bien attestés au Moyen Âge. À l’origine, ce sont peut-être des tenures de défrichements, comparables à celles qui sont repérables aujourd’hui par un anthroponyme précédé de « chez », et qu’on trouve fréquemment dans d’autres parties de la Savoie. Mais dans ce cas précis, le « Les » montre bien, nous semble-t-il, qu’après l’installation d’une première famille de colons, les descendants ont proliféré jusqu’à faire d’une tenure primitive un véritable village. »\footcite[247]{Carrier2001VieMontagnardeFaucigny}
\end{quote}

Dans son étude sur le Faucigny médiéval, Nicolas Carrier propose que certains mas et communaux puissent trouver leur origine dans d'anciennes exploitations familiales demeurées longtemps indivises. Il écrit ainsi :

\begin{quote}
	« D’où la tentation de penser que les communaux de hameaux, qui sont si nombreux aux XIV\textsuperscript{e} et XV\textsuperscript{e} siècles, sont, pour une bonne partie d’entre eux, d’anciennes propriétés familiales restées indéfiniment indivises. [...] Il semble donc bien que les communaux de prises et de mas correspondent à la part des terres de ces exploitations qui n’a jamais été partagée. »\footcite[247]{Carrier2001VieMontagnardeFaucigny}
\end{quote}

Plus loin, l'auteur souligne également que de nombreuses communautés de hameaux pourraient avoir eu une origine lignagère :

\begin{quote}
	« De nombreux hameaux faucignerans ont donc une origine familiale. Ils sont nés et se sont développés parce que des cognats continuaient à voisiner après avoir quitté la cohabitation, et gardaient sur de nombreuses générations une partie de leurs biens en communs, ainsi devenus communia de hameaux. »\footcite[247-248]{Carrier2001VieMontagnardeFaucigny}
\end{quote}

Une telle hypothèse soulève naturellement la question de savoir si des traces de ces communautés familiales anciennes demeurent perceptibles dans le Beaufortain du XVIII\textsuperscript{e} siècle. Si certains mas ont effectivement pour origine des groupes de parenté ayant reçu ou défriché collectivement un territoire au Moyen Âge, on peut s'attendre à y retrouver, plusieurs siècles plus tard, une concentration anormalement forte de certains patronymes.

La Mappe Sarde offre la possibilité d'explorer cette question à travers les propriétaires déclarés en 1730. L'objectif n'est pas de démontrer l'existence de telles communautés lignagères, mais de vérifier si leur éventuelle existence médiévale demeure encore perceptible dans la structure de la propriété foncière observée au XVIII\textsuperscript{e} siècle.

\begin{table}[H]
	\centering
	
\begin{tabular}{lrrlr}
\toprule
Patronyme & Nb_prop & Surface & Mas & Surface_Mas\\
\midrule
BLANC & 18 & 1203.26 &  & \\
 &  &  & Roniaux & 291.48\\
 &  &  & Traicol & 249.25\\
 &  &  & Lay & 180.02\\
 &  &  & Couecle & 101.6\\
\addlinespace
 &  &  & Borne Et Meralliet & 78.1\\
 &  &  & Val De Treicol & 63.12\\
 &  &  & Planay Voisin De La Grande Combe & 35.23\\
 &  &  & Roselend & 34.21\\
 &  &  & Voniaret & 30.72\\
\addlinespace
DOIX & 22 & 610.05 &  & \\
 &  &  & Berge & 115.72\\
 &  &  & Charmette & 62.54\\
 &  &  & Val & 60.17\\
 &  &  & Bettiere & 52.73\\
\addlinespace
 &  &  & Roniaux & 38.1\\
 &  &  & Gitte & 33.26\\
BOCHET & 29 & 552.13 &  & \\
 &  &  & Lay & 226.79\\
 &  &  & Grande Combe & 177.17\\
\addlinespace
 &  &  & Planay Voisin De La Grande Combe & 38.61\\
GACHET & 26 & 384.89 &  & \\
 &  &  & Gitte & 162.48\\
 &  &  & Sceause & 77.89\\
 &  &  & Roselend & 42.67\\
\addlinespace
 &  &  & Val De Treicol & 42.14\\
VIBERT & 32 & 247.56 &  & \\
 &  &  & Val & 49.57\\
 &  &  & Val De Treicol & 44.37\\
VIALLET & 39 & 231.21 &  & \\
\addlinespace
 &  &  & Cres Denset & 64.31\\
 &  &  & Bouchères & 30.03\\
NEIROUD & 15 & 180.28 &  & \\
 &  &  & Borne Et Meralliet & 58.25\\
 &  &  & Perrosan & 57.43\\
\addlinespace
MOLLIET & 18 & 87.99 &  & \\
 &  &  & Lay & 56.03\\
JOLY & 18 & 78.56 &  & \\
 &  &  & Bettiere & 33.34\\
\bottomrule
\end{tabular}

	\caption{Principaux patronymes et répartition par mas}
	\label{tab:PatronymesMas}
\end{table}

Le tableau présente les patronymes associés à plus de quatorze propriétaires distincts dans l'ensemble de la paroisse (hors « Institution »). Pour chaque patronyme sont indiqués le nombre total de propriétaires correspondants et la surface totale détenue. Ne sont ensuite affichés que les mas dans lesquels ce patronyme possède au moins 30 hectares, avec la surface correspondante.

Le tableau \ref{tab:PatronymesMas} présente les principaux patronymes de la paroisse et leur répartition foncière entre les différents mas. Si certaines familles semblent particulièrement associées à certains secteurs — les Blanc à Traicol et à Lay, les Doix à Berge ou Charmette, les Viallet à Cres Denset ou Bouchères — aucune n'apparaît véritablement cantonnée à un territoire précis.

Au contraire, les principaux patronymes sont présents dans de nombreux mas distincts, parfois éloignés les uns des autres. Cette dispersion géographique traduit une circulation ancienne des patrimoines fonciers, résultant probablement des successions, des alliances matrimoniales, des ventes, des échanges et des recompositions d'indivisions.

Cette observation est particulièrement intéressante au regard des hypothèses formulées pour d'autres régions alpines, où certains auteurs ont proposé une origine lignagère des mas et des communaux. Si certains territoires ont pu être initialement concédés à un groupe familial ou à une communauté de parents, la situation observée à Beaufort en 1730 ne permet plus d'en retrouver clairement la trace à travers la seule répartition des patronymes.

\begin{table}[H]
	\centering
	
\begin{tabular}{l}
\toprule
\textbf{Gitte} : 1889.6 ha ; 13 patronymes ; 20 propriétaires\\
\hspace{1em}GACHET (1 prop.) : 162.5 ha\\
\hspace{1em}BOUCHAGE (3 prop.) : 143 ha\\
\hspace{1em}DOIX (2 prop.) : 33.3 ha\\
\\
\addlinespace
\textbf{Lay} : 1199.5 ha ; 8 patronymes ; 11 propriétaires\\
\hspace{1em}BOCHET (1 prop.) : 226.8 ha\\
\hspace{1em}LANCHE (2 prop.) : 221.3 ha\\
\hspace{1em}BLANC (3 prop.) : 180 ha\\
\hspace{1em}CORNU (1 prop.) : 151.2 ha\\
\addlinespace
\hspace{1em}MOLLIET (1 prop.) : 56 ha\\
\\
\textbf{Ladray} : 891.8 ha ; 23 patronymes ; 65 propriétaires\\
\hspace{1em}PERRIER (11 prop.) : 56.8 ha\\
\\
\addlinespace
\textbf{Roselend} : 457.3 ha ; 25 patronymes ; 49 propriétaires\\
\hspace{1em}BOUCHAGE (7 prop.) : 78.7 ha\\
\hspace{1em}GACHET (3 prop.) : 42.7 ha\\
\hspace{1em}BLANC (3 prop.) : 34.2 ha\\
\hspace{1em}LANCHE (3 prop.) : 25.5 ha\\
\addlinespace
\\
\textbf{Traverse} : 248 ha ; 21 patronymes ; 43 propriétaires\\
\hspace{1em}DUC (13 prop.) : 27.4 ha\\
\\
\textbf{Val De Treicol} : 213.6 ha ; 6 patronymes ; 10 propriétaires\\
\addlinespace
\hspace{1em}BLANC (3 prop.) : 63.1 ha\\
\hspace{1em}VIBERT (2 prop.) : 44.4 ha\\
\hspace{1em}GACHET (1 prop.) : 42.1 ha\\
\\
\textbf{Val} : 187.4 ha ; 5 patronymes ; 6 propriétaires\\
\addlinespace
\hspace{1em}DOIX (1 prop.) : 60.2 ha\\
\hspace{1em}VIBERT (2 prop.) : 49.6 ha\\
\\
\textbf{Bouchères} : 176.5 ha ; 11 patronymes ; 19 propriétaires\\
\hspace{1em}VIALLET (8 prop.) : 30 ha\\
\addlinespace
\\
\textbf{Cres Denset} : 128.6 ha ; 4 patronymes ; 9 propriétaires\\
\hspace{1em}VIALLET (5 prop.) : 64.3 ha\\
\\
\textbf{Charmette} : 115.5 ha ; 7 patronymes ; 10 propriétaires\\
\addlinespace
\hspace{1em}DOIX (3 prop.) : 62.5 ha\\
\\
\textbf{Leila} : 104.2 ha ; 6 patronymes ; 13 propriétaires\\
\hspace{1em}VIALLET (8 prop.) : 28.9 ha\\
\\
\bottomrule
\end{tabular}

	\caption{Grands mas et principaux patronymes associés}
	\label{tab:MasPatronymes}
\end{table}

Le tableau recense les mas de plus de 100 hectares et indique, pour chacun d'eux, sa surface totale, le nombre de patronymes et le nombre de propriétaires distincts (hors « Institution »). Seuls sont retenus les patronymes comptant au moins dix propriétaires dans l'ensemble de la paroisse et possédant au moins 20 hectares dans le mas considéré. Pour chaque patronyme est également indiqué le nombre de propriétaires distincts concernés dans le mas.

Le tableau \ref{tab:MasPatronymes} aborde la question sous l'angle inverse du précédent en partant des plus grands mas de la paroisse. Si l'hypothèse d'une origine familiale des mas demeurait encore visible au XVIII\textsuperscript{e} siècle, on pourrait s'attendre à observer, dans les grands territoires d'altitude, une forte concentration de propriétaires portant un même patronyme.

Or les résultats obtenus montrent généralement une situation beaucoup plus complexe. Les grands mas les plus étendus, tels que la Gitte, Roselend, Ladray ou Traverse, rassemblent un nombre élevé de patronymes et de propriétaires distincts. Même lorsque certains noms apparaissent plus fréquemment que d'autres, ils ne représentent qu'une fraction limitée des ayants droit.

Quelques exceptions existent, notamment à Cres Denset, Charmette ou dans le Val de Treicol, où certains patronymes semblent plus fortement représentés. Toutefois, ces cas demeurent isolés et ne suffisent pas à faire apparaître une organisation générale fondée sur des groupes familiaux identifiables.

Ainsi, l'étude des patronymes ne permet pas de mettre en évidence, dans le Beaufort de 1730, la survivance manifeste de communautés lignagères associées à un mas particulier. Si de telles structures ont existé au Moyen Âge, plusieurs siècles de transmissions successorales, de mariages, de ventes et de recompositions foncières semblent avoir largement effacé leur lisibilité. Les grands mas apparaissent désormais comme des ensembles fonciers complexes associant de nombreux lignages plutôt que comme les territoires propres à une famille ou à un groupe de parenté.

\section{Les Communs}

\begin{quote}
	« À partir du XIV\textsuperscript{e} siècle, là et seulement là où est passé le mouvement des franchises et des albergements, les eaux, les forêts, les pâturages exploités par les paysans sont désormais reconnus comme des communia, c’est-à-dire des communs. Les communautés en détiennent la propriété utile ce qui inclut le droit d’accès, plus ou moins exclusif, et le droit de police, c’est-à-dire la faculté de réglementer l’exploitation de ces espaces et de réprimer le manquement à ces règles. En Savoie, la communauté paroissiale apparaît comme le cadre naturel de la possession et de la gestion de ces droits. Ailleurs, les bois et les alpages peuvent relever de plusieurs niveaux de communauté : celui de la vallée, celui de la paroisse et celui du village. En Savoie, les communautés de vallée sont l’exception et, d’une façon générale, on n’y trouve pas, ni d’ailleurs dans aucune autre partie des Alpes françaises, de biens collectifs relevant de l’ensemble d’une vallée. Cette situation existe pourtant depuis le Moyen Âge dans les Alpes orientales italiennes et autrichiennes ou encore dans les Pyrénées occidentales, par exemple dans la vallée d’Ossau. En Savoie, la gestion de l’espace commun est l’affaire de deux institutions qui peuvent agir de manière indépendante ou complémentaire, à savoir la communauté d’habitants et ce que l’on appelle ici les pareries ou, plus couramment, les consortages. […]
	
	En Savoie, la communitas prend donc la paroisse pour cadre exclusif. À tel point qu’au XIV\textsuperscript{e} et au XV\textsuperscript{e} siècle, une communauté de quartier ou de village particulièrement affirmée, va tout faire pour obtenir le statut paroissial, exigeant de son évêque l’érection de la modeste chapelle de hameau en église de plein droit, ou tout au moins en église filiale, dotée d’un cimetière et de fonds baptismaux. La Giettaz, jusque-là quartier de la paroisse de Flumet, devient paroisse annexe en 1390 ; Morzine, quartier de Saint-Jean-d’Aulps, en 1499 ; Morillon »
	\footcite[85]{Mouthon2016}
\end{quote}

\begin{quote}
	« Que l’on n’imagine pas pour autant les communautés de paroisses comme de véritables municipalités rurales pourvues d’un exécutif doté de prérogatives étendues. En cette fin du Moyen Âge, le pouvoir communautaire réside bien davantage dans l’assemblée des chefs de famille que dans les procureurs syndics ou les prud’hommes que celle-ci désigne occasionnellement et auxquels elle n’octroie qu’un mandat limité. Entre la fin du XIII\textsuperscript{e} et le début du XVI\textsuperscript{e} siècle cependant, à l’imitation des villes, les communautés se dotent peu à peu d’un ou deux syndics permanents, choisis pour un an, plutôt dans les familles de notables, et dotés de pouvoirs plus généraux. Ce n’est qu’avec les réformes administratives des années 1560, et la généralisation des syndics, que les communautés paroissiales savoyardes deviennent officiellement des communes. Or ce sont ces collectivités, à l’organisation encore incertaine, qui, à la fin du Moyen Âge, se sont vu reconnaître ou confirmer le contrôle des « choses de personne » (\textit{res nullius}) situées sur leur territoire. Lors de l’albergement ou par l’octroi des franchises, le seigneur cède la propriété utile sur les cours d’eau, les forêts, les alpages. Il conserve seulement le droit dit « éminent » ou encore la « directe » dont le cens en argent payé chaque année par la communauté constitue la reconnaissance à la fois symbolique et monétaire. »
	\footcite[86]{Mouthon2016}
\end{quote}


\section{Les propriétaires}
	
	Avant de commercer l’exploration de la propriété des alpages, nous allons essayer de comprendre la propriété dans l’ensemble de Beaufort :
	
	\begin{table}[htbp]
		\centering
		\includegraphics[width=0.5\textwidth]{tables/1 Propriétaires par Statut.JPG}
		\tabcaptionindex{tab:PropStatut}{répartition du nombre de propriétaire par statut}
	\end{table}
	tableau : répartition du nombre de propriétaires par statut
	
	
	a) Propriétaires déclarés « Ecclésiastique »
	
	Les propriétés ecclésiastiques sont quasi inexistantes (8 ha) et les parcelles sont en grande majorité du bâti (église, chapelles…) et des structures attenante (jardin, cour, cimetière…). Aucun grand propriétaire terrien ecclésiastique, pas de communauté religieuse, ni de biens appartenant à un monastère ou au diocèse de Moutier.

	La liste des 22 propriétaires ecclésiastiques est la suivante :
	
	La Cure de Beaufort
	Chapelle du Bersend
	Chapelle des Pénitents du St Sacrement
	Chapelle d'Areche
	Chapelle de N. D. de pitié
	Chapelle de la Portelaz
	Chapelle de la Pierre
	Chapelle de Roselend
	Chapelle à la Dray
	Chapelle de la Frasse
	Chapelle des Curtillets
	Chapelle des Houtards
	Chapelle des villes dessus
	Chapelle des villes dessous
	Chapelle de la Gitte
	Chapelle des Chornets
	Chapelle du Mont
	et Les Rvds Cordelliers de Ste Marie Egyptienne de Chambéry
	
	Les Cordeliers de Chambéry sont la seule communauté religieuse externe et elle ne possède que deux masures avec pré et jardin.
	La présence religieuse est importante au vu du nombre de chapelles mais cela reste cohérent par rapport à la superficie de la commune, la dispersion de l’habitat et les difficultés de se rendre à l’église de Saint-Maxime en hiver.
	
	En revanche 8 prêtres sont propriétaires (seuls et parfois aussi en indivision), mais sont déclarés comme "communier". Cela ne signifie pas forcément qu'ils exercent sur la paroisse, mais on sait que les habitants de Saint-Maxime ont obtenu le droit de choisir leurs prêtres.
	
	Ce qui est le cas du "Rvd LANCHE curé de Beaufort" élu par les paroissiens en 1717.
	
	
	b) Propriétaires déclarés « Noble »
	
	Les propriétaires nobles sont aussi très anecdotiques :
	
	\begin{table}[htbp]
		\centering
		\includegraphics[width=0.5\textwidth]{tables/1 Prop Nobles.JPG}
		\tabcaptionindex{tab:PropNoble}{Liste des propriétaires nobles}
	\end{table}
	

	Le Marquis de Beaufort (dont le premier titre est Marquis de Fleury) ne possède qu’un hectare alors qu’il détient la plupart des droits féodaux sur le mandement de Beaufort.
	
	Gaspard LAURENT(200ha) et le Marquis de Cordon (99ha) sont des propriétaires importants, mais il ne semble pas qu’il s’agisse de terres associées à leur historique féodal. 

	Le Marquis de Cordon est le seul noble possédant des terres de manière significative ainsi que des droits féodaux sur la commune. Cependant Hélène Viallet signale qu’il est propriétaire d’un alpage de 60 vaches vendu en 1780 \footcite[102]{HViallet1993} mais elle n’en fait pas état dans le cadastre de 1607 ni dans celui de 1645.
	\footnote{François Joseph Sallier de La Tour (1706 - 1799), Marquis de Cordon ; ADS -208 Edépôt 286 Fief de la dame de Saint-Sulpice et du marquis de Cordon. – Etat estimatif de la rente et fief que la comtesse d'Ugine, épouse du seigneur de Saint-Sulpice, a éteint et aboli en faveur des communautés du Villard et Saint-Maxime de Beaufort (1773). Répartition du prix du fief entre les communautés, quittance (1779-1780). Etat de la rente féodale appartenant au marquis de Cordon (1782)} : .
	
	Nous sommes donc loin d’une propriété issue de la féodalité tardive, avec de grands propriétaires fonciers nobles desquels dépendent de nombreux fermiers.
	
	En revanche, concernant le marquis de Cordon et Gaspard Laurent, il est possible qu’il s’agisse d’acquisitions de type  investissement dans les alpages, en parallèle de celles des notaires suite à l’introduction de la fabrication du gruyère dans le Beaufortain.
	
	\begin{quote}
		« Au milieu du XVI\textsuperscript{e} siècle, un quart des nobles tient du bétail, contre 80 \% pour l’ensemble de la population. Ce taux n’est pas négligeable et il atteint 40 \% dans les paroisses rurales. On a même recensé huit feux nobiliaires déclarés pauvres ou misérables. Pauvre ou pas, un noble conserve son rang et ses privilèges, même lorsqu’il travaille la terre ou élève quelques bêtes ; ce n’est pas une activité de dérogeance, à la différence du notariat, du commerce et des arts mécaniques. [...] Quant à la fonction de châtelain, nous l’avons vu, elle a complètement perdu de son prestige au milieu du XVI\textsuperscript{e} siècle et on ne compte que cinq châtelains nobles dont un seul participant aux dénombrements, celui de Rumilly. »
		\footcite[142]{Gachet2011GabelleSel1561}
	\end{quote}
	

	c) Propriétaires déclarés « Bourgeois » ou « Forains »
	
	
	
	\begin{table}[htbp]
		\centering
		\includegraphics[width=0.5\textwidth]{tables/1 Prop Bourgeois et Forains.JPG}
		\tabcaptionindex{tab:PropBourg}{Liste des propriétaires "bourgeois" et "forains"}
	\end{table}
	
	
	Nous y trouvons un notaire et deux avocats.
	Nous investiguerons sur le plus gros propriétaire, Claude BLANC, bourgeois de Montmeillant (aujourd’hui Montmelian à coté de Chambéry), il est possible et probable qu’il soit originaire de Beaufort, mais qu’il se soit installé proche de la capitale administrative qu'est Chambéry. Nous y reviendrons lorsque nous étudierons le cas de la famille BLANC.
	Nous avons donc 4 particuliers hors des communiers qui sont de gros propriétaires : deux nobles et deux bourgeois totalisants à eux quatre :1 179ha (soit \texttt{8\%} de la surface de la commune)
	
	
	
	c) Propriétaires déclarés « Communaux » et les « Communiers»
	
	Il reste deux catégories : les communaux et les communiers à part presque égale en termes de superficie.
	
	
	Tableau : Surface en ha par statut de propriétaires et degré de bonté
	
	
	La propriété des surfaces incultes (degré de bonté 0) est répartie entre les statuts en proportion des totaux. Les communiers possèdent la plus grande part de terres incultes xxxxx???? à vérifierXXXX.
	
	
	On remarquera par la même occasion que l’administration ducale a considéré que la commune de Beaufort était très pourvue en surface d’excellente qualité : \texttt{68\%} de la surface a obtenu le degré de bonté maximum de 3.
	
	tableau : répartition des surfaces entre statut des propriétaires et Occupation du Sol

	Si en global, les surfaces des communaux et celles possédées par les communiers étaient du même ordre de grandeur. Il en est différemment lorsque nous regardons le détail par l’utilisation du sol. Ainsi les communaux détiennent \texttt{78\%} des bois et \texttt{58\%} des rocs. Les communiers possèdent \texttt{52\%} des pâturages (et en ajoutant les « Autres » dont principalement deux bourgeois et deux nobles, nous atteignons \texttt{63\%} des surfaces en pâturage.
	
	Répartition des propriétaires « Communier » en fonction de la superficie possédée : 
	Les communiers possèdent en moyenne 7.16ha (avec une moyenne de 19 parcelles par propriétaire)
	
	tableau d’analyse des exploitations des communiers par quartile.

	Nous constatons que dans le quartile supérieur, \texttt{25\%} des propriétaires possèdent \texttt{85\%} des surfaces avec une variance extrême. Puisque l’exploitation la plus petite ne fait que 4.38ha (ce qui permet de faire paître peut-être 1 ou 2 vaches seulement et à condition d’avoir accès aux pâturages communaux). Alors que la plus grosse exploitation dispose de 265ha.
	
\section{Les propriétaires en fonction de l’altitude}
	
	En premier lieu, la superficie des terrains au-dessus de 1600m représente 67% de la surface totale de la commune (9 367ha sur 13965ha)
	En dessous de 1600m d’altitude, nous avons 697 communiers propriétaires de bois, broussailles, pâturage ou pré (dont certains peuvent avoir aussi des terrains au-dessus de 1600m). mais toutes les propriétés sont petites (< 5ha seuls 49 propriétaires ont plus de 5ha dont seulement 14 ont plus de 10ha ; la plus grande propriété a 28ha)
	En revanche au-dessus de 1600m nous avons de grandes propriétés appartenant à des communiers. 155 communiers possèdent des terres à bois, broussailles, pâturages ou prés au-dessus de 1600m dont 50 ont plus de 10ha.
	
	Comment définir la notion d’alpages (ou de grand territoire au-dessus de 1600m) ?
	Il y a plusieurs réponses possibles, sachant que la taille moyenne d’une exploitation agricole en Savoie au XVIIIe n’est pas forcément significative. En revanche, nous nous intéressons ici aux changements liés à l’introduction du gruyère suisse dans le Beaufortain. Et pour fabriquer ce fromage selon les règles, il faut suffisamment de vaches pour que le fruit d’une traite permette de fabriquer un fromage. Aujourd’hui pour respecter les règles de l’AOC, il faut environ 50 vaches. Au XVIIIe siècle, la taille des meules de Beaufort n’était pas figée, les meules étaient plus petites (40kg aujourd’hui, entre 20 et 30kg à l’époque ???) mais dans l’autre sens, la productivité des vaches du XIIIe était sans doute inférieure à celle d’aujourd’hui. Donc nous pourrions chercher les pâturages dont la superficie en altitude permettrait de nourrir pendant les trois mois d’été une quarantaine de vaches.
	Mais l’accès aux pâturages communaux permet à n’importe quel propriétaire de la vallée d’y faire paître son troupeau. Nous savons par les plaintes concernant les abus d’utilisation des communaux, que les propriétaires de parcelles privées situées à proximité des communaux sont les principaux bénéficiaires de ceux-ci. Ils peuvent faire monter leur troupeau avant la date de début générale de la remue, et de même redescendre leurs troupeaux après la date de fin.
	Est-ce que 10ha de pâturage en altitude (alpage) possédés par un communier permettent de le qualifier « grand propriétaire » ? 
	
	
	 tableau : répartition des surfaces >1600m entre statut des propriétaires et Occupation du Sol
	Au-dessus de 1600m, il n’y a plus de champs, presque plus de pré et proportionnellement moins de broussailles.
	Il ne reste que des Pâturage, des forêts et du rocher 
	Ci-dessous figure la liste des propriétaires (de tous statuts) qui possèdes des pâturages de plus de 10ha au-dessus de 1600m :
	
	tableau des « grands » propriétaires d’alpage : surface totale possédée supérieure à 10ha et située au-dessus de 1600m
	Il y a 52 entrées :
	-	2 Bourgeois
	-	2 Nobles
	-	1 Forain (marchand de la commune mitoyenne de Granier 
	-	0 bien ecclésiastique
	-	2 communaux : la communauté de Saint-Maxime mais aussi, déclaré comme communaux un pâturage appartenant aux « communiers de la commune de Granier »
	-	45 communiers de Saint-Maxime.
	
	Nous remarquons la présence de propriétaires avec des noms de famille communs : les patronymes BLANC et DOIX sont les plus représentés, nous verrons que certains sont de la même famille.
	Nous pouvons remarquer que ces « grands » propriétaires possède leur « domaine » en majorité en altitude :
	28 propriétaires sur 52 ont plus de \texttt{90\%} (dont 11 d’entre eux en ont \texttt{100\%})
	40 propriétaires sur 52 ont plus de \texttt{80\%} de leur patrimoine au-dessus de 1600m (28 plus de \texttt{90\%} et même 11 sont à \texttt{100\%} !)
	Nous voyons deux interprétations possibles : 
	-	Une première interprétation possible est que certains propriétaires sont des « entrepreneurs », ils ne sont pas exploitants ni paysans mais : bourgeois, forain, marquis, prêtres ceux qui acquièrent des alpages avec l’objectif de faire fabriquer du « gruyère » façon Beaufort, ils ne s’intéressent qu’à l’investissement dans les alpages et ils dispersent le troupeau pour l’hiver chez de petits propriétaires dans la vallée. 
	-	Une autre explication provient du fait que les techniques agricoles du XVIIIe siècle ne permettent pas de grandes étables afin d’y loger de grands troupeaux, il faut faucher le foin à la faux (donc manuellement) il faut transporter ce foin à dos de mulet ou à dos d’homme. Les étables permettent d’hiverner quelques vaches seulement par chalet. L’accès par la propriété aux alpages d’altitude permets à une catégorie « supérieur » d’éleveur de rassembler en « pension » pour l’été les vaches de tous les petits propriétaires de la vallée.
	À peu près toutes les combinaisons possibles ont été imaginées par les propriétaires des alpages pour contractualiser l’inalpage et la répartition des fruits (le fromage) les naissances des veaux, la propriété d’une vache être morcelée (un tel possède ¼ de vache). Le propriétaire d’alpage pourrait en théorie ne posséder aucune vache en propre, mais inalper uniquement des vaches de ses « clients). De l’autre côté, le petit propriétaire d’une surface de fauche en vallée peut loger un mixte de vaches lui appartenant et de vaches appartenant à d’autres. 
	Ce qu’il est important de préciser est que ce mode de gestion des alpages est bien antérieur à la fabrication du « Beaufort » mais l’introduction de Suisse de ce fromage à pâte cuite qui se conserve très longtemps et donc qui s’exporte, a attiré des nouveaux acteurs, investisseurs et commerçant.
	
	
	
\section{Qui possède les alpages ?}

		\begin{quote}
		« À l’inverse, Hélène Viallet cite le cas d’une part d’alpage du Beaufortain devenue grande montagne au milieu du XV\textsuperscript{e} siècle. En 1448, l’alpage du Bouchoz au-dessus du village de mi-saison de Roselend (englouti depuis 1960 sous un lac de barrage) est tenu par trente-cinq consorts qui se partagent dix-sept parts. La reconnaissance fait référence à Pierre Vial et audit Neyroz, son frère, qui, à l’origine (au XIV\textsuperscript{e} siècle ?), avaient reçu ce qui était alors une part de petite montagne, dans le cadre d’un « antique albergement ». La règle successorale imposant partout le partage des biens entre les fils (les filles étant en principe dotées et mariées), le Bouchoz, qui n’était au départ qu’une portion d’une petite montagne, sans doute assez vaste, est devenue une « petite grande montagne » par la volonté des héritiers de rester dans l’indivision. Ce type d’albergement d’un pâturage de faible superficie à une seule personne ou une fratrie semble avoir été à l’origine de plusieurs montagnes à consortage, toujours par le jeu des partages d’héritage. »
		\footcite[120]{Mouthon2019MontagnesMedievales}
	\end{quote}
	
	\begin{quote}
		« Le caractère cessible des parts d’alpage ouvre également la porte à leur rachat et à leur accumulation par les plus riches des éleveurs. C’est ce qui se produit à partir du XVI\textsuperscript{e} et surtout du XVII\textsuperscript{e} siècle. Cette tendance, qui mène à la constitution des montagnes particulières ou privées, n’est pas pour autant générale, tant s’en faut. Elle caractérise uniquement une partie des montagnes à consortages et pas les grandes montagnes communales. Elle semble également concerner en premier lieu le Beaufortain et le Haut-Faucigny, mais pas la Tarentaise. En 1552, une supplique des habitants du mandement de Beaufort explique de façon très claire que les plus riches d’entre eux accaparent les pâturages communs aux dépens des plus pauvres : « icelles communautés se gaudissent et jouyssent par gens du pays riches ayant puissance de bestail de sorte que le povre menu peuple d’icelles est quasiment frustré ». »
		\footcite[125]{Mouthon2019MontagnesMedievales}
	\end{quote}
	
	
	
	\begin{quote}
		« Le mode d'exploitation le plus fréquent pour une grande montagne est donc le fermage, dénommé en Savoie acensement ou admodiation. C'est toujours un acte écrit passé devant notaire en raison du montant élevé de la censé. [...] Le 22 décembre 1755, Joseph Blanc ascense sa 	montagne de Conchettes à Joseph Chevallier-Joly, sous la censé de 44 quintaux 	de gruyère à son choix (Archives Blanc)»
		\footcite[113]{HViallet1993}
	\end{quote}
	
	
	\begin{quote}
		« Le livre de raison de Me Joseph Blanc(36), secrétaire insinuateur et plus riche propriétaire en bétail de la vallée nous donne des indications sur le lieu d'hiverné de ses vaches : en 1771, il a placé cent dix-huit vaches ; quatre à Palud, neuf à Saint-Sigismond, trois à l'Hôpital, une à Villard, trois à Hauteluce et quatre-vingt quatorze à La Pierre (dans quarante foyers). »
		\footcite[58]{HViallet1993}
	\end{quote}
\section{Analyse}

\section{Cas Particulier de la famille Blanc}
	-	histoire de notable notaire insinuateur
	-	Mappe Sarde
	-	Mappe Blanc
	

	
	\cartefigure[0.3\textwidth]
	{figures/4 extrait 1.jpg}
	{fig:ParcellesJABlanc}
	{Extrait des parcelles du Rdv Jean-Antoine BLANC}
	
	\cartefigure[0.3\textwidth]
	{figures/4 extrait 2.jpg}
	{fig:ParcelleFBlanc}
	{Extrait des parcelles de Françoise CHRISTINé veuve de Michel BLANC et possesseurs voisins }
	
	\cartefigure[0.3\textwidth]
	{figures/4 extrait 3.jpg}
	{fig:ParcellesChapelleBlanc}
	{Extrait des parcelles qui ont été rayées attribuées à la Mappe des Chapelles }
		
	\cartefigure[0.3\textwidth]
	{figures/4 Genealogie_Blanc.jpg}
	{fig:GenealogieBlanc}
	{Arbre généalogique de la famille Blanc (XVII\textsuperscript{e}--XIX\textsuperscript{e} siècles)}
	
	
\chapter{Litiges des chapelles}	
	
	
	\begin{quote}
		« La vallée de Treicol est un secteur âprement disputé par Aime et les Chapelles, contre Saint-Maxime. Elle fait l'objet d'un procès pluri-séculaire, puisqu'il débute en 1599 par le refus de Mc Jean-François Dunand de payer la taille de sa montagne de Presset à Beaufort. L'affaire s'envenime, car de nombreux particuliers possèdent des alpages à Treicol. Pour être limitrophe de la vallée et augmenter ses prétentions, la communauté des Chapelles va jusqu'à acheter en 1732 une longue et étroite bande de rochers depuis la Pierre-Menta jusqu'à son territoire. Portée devant toutes les instances possibles, l'affaire n'est définitivement tranchée qu'en 1907, par un arrêt du Conseil d'Etat qui attribue la vallée à Beaufort. »
		\footcite[98]{HViallet1993}
	\end{quote}
	
	
\chapter{18e Evolution}


\begin{quote}
	« La vente des biens communaux est soumise à un certain formalisme. Des experts et géomètres choisissent les parcelles, qui font l’objet d’une publication et d’un affichage au banc du droit de la ville. La vente est publique, aux enchères. Celle qui a lieu en 1773 donne lieu à quelques troubles : les paysans trop pauvres pour acquérir des parcelles manifestent et le conseil doit s’assurer le concours d’un brigadier et de deux soldats les jours d’enchère. En mars 1778, d’autres biens communs sont vendus. Ce sont les gros montagnards qui profitent surtout des ventes. Les plus gros achats sont effectués par Blanc Joseph, secrétaire insinuateur, Chevallier-Joly Jean, notaire, Bouchage André, Blanc Jean-Baptiste, Doix Claude et le révérend Molliet Joseph. Quelques petits propriétaires acquièrent également de modestes parcelles en indivision. Environ soixante-dix familles profitent de ces ventes. Saint-Maxime, pour une cession de 1\,148 hectares, obtient une somme de 58\,868 livres et 11 sols ; déduction faite des frais dus à l’achat d’une maison commune et aux différentes vacations, il reste à la communauté 54\,091 livres et 19 sols. (A.D. de la Savoie, C 4994 — Communaux de Saint-Maxime de Beaufort vendus pour le payement du prix des affranchissements en 1782)»
	\footcite[127]{BordeVibertGuigue1971VieRuraleBeaufort}
\end{quote}


\section{Bruchet1908Abolitions}
\begin{quote}
	« Toutefois, pour comprendre l’importance du sacrifice consenti par certaines communautés pour se libérer, il faut tenir compte de la proportion entre le montant des affranchissements et leurs revenus. La taille foncière calculée à raison du 1/5 moins 1/76 de revenu donne une indication utile. On voit par exemple que S. Maxime-de-Beaufort, qui eut à supporter 44\,610 l. d’affranchissements, put rapidement rembourser les seigneurs parce que c’était une localité riche dont la taille s’élevait à 9\,862 livres : l’affranchissement de cette localité représentait une somme variant entre quatre et cinq fois son impôt foncier alors que, dans d’autres pays de la Savoie-propre, la proportion atteignait vingt fois, cinquante fois et quatre-vingts fois l’impôt foncier. Ces énormes sacrifices consentis par les populations prouvent mieux que tout autre argument la nécessité de la réforme poursuivie par la maison de Savoie. »
	\footcite[LXXXI]{Bruchet1908Abolition}
\end{quote}

\begin{quote}
	« S. Maxime-de-Beaufort (Savoie, aujourd’hui Beaufort). — Vicardel marquis de 
	Fleury 36\,705 l. ; Chapelle de S. François 305 l. ; du Belex 260 l. ; Chapelle 
	des 10\,000 Vierges à Aime 305 l. ; Chamiol-Maitral 3\,250 l. ; Vuichard de 
	S. Réal et de Burgarel de Beaufort 8\,370 l. ; Chevillard ép. Morand de 
	S. Sulpice 620 l. ; Bucher de S. Guillaume veuve Bertrand de Chamoussel 
	1\,220 l. ; Archevêché de Moûtiers 4\,500 l. En 1791, cette communauté avait 
	entièrement payé son affranchissement sauf 565 l. Total : 55\,515 l., dont 
	44\,610 l. supportées par la seule communauté de S. Maxime-de-Beaufort qui 
	paya complètement cette somme. 
	Le marquis de Fleury se réservait toutefois « la recherche des cristaux, le droit d'avoir un banc placé ainsi qu'il convient à la qualité du seigneur dudit lieu dans l'église paroissiale du bourg de Saint-Maxime..., comme encore le droit de faire placer des patibulaires entre ses armoiries.	\footcite[426]{Bruchet1908Abolition}
\end{quote}

\section{Rachat des droits féodaux-contexte}
\section{Cas du Beaufortain}
\begin{quote}
	« Le conseil de Saint-Maxime avait décidé dès 1766 le rachat du fief : « pour l'utilité générale du public, rien n'est préférable à la liberté de pouvoir sans obstacle disposer des biens à chacun départis, et de les affranchir de cette foule de devoirs seigneuriaux qui ne laissent à ce pays de montagne où les biens sont d'un culture difficile et d'un petit rapport, pour toute récompense et dédommagement au propriétaire, que la peine de les cultiver ».
	
	Cette opération est financée par la vente d'une partie des immenses communaux, bien que cela suscite l'opposition de certains communiers<56). Le 15 mai 1770, le contrat d'affranchissement envers le marquis de Fleury est enregistré par la Chambre des comptes de Turin. La communauté a racheté le principal fief au prix de 36.705 £ 4s, et cela six mois avant l'Edit d'affranchissement du 19 décembre 1771. La politique de réduction des privilèges nobiliaires permettait aux communautés d'imposer à leur seigneur le rachat des censés, servis, laods et autres droits. Il suffisait de l'avis favorable des deux tiers des communiers réunis en assemblée générale.
	
	Cette réussite engagea le conseil à poursuivre les opérations de vente des communaux en vue du rachat des autres fiefs, réalisé en 1773. Le prix total des affranchissements, 47.384 livres 9 sous 5 deniers, était entièrement payé en 1781.
	
	De 1773 à 1778, il avait fallu vendre 15 \% des communaux soit 1 026 hectares. A cette occasion, certains notables avaient su tirer un parti très avantageux des réformes politiques édictées par le souverain, en rachetant les meilleures portions des communaux aliénés. Si la vente des biens communaux n'était pas allée dans le sens de leurs intérêts personnels, les administrateurs de la communauté n'auraient pas mené à bien les opérations longues et difficiles du fait des oppositions. Le rachat de la liberté, présenté d'une façon très démagogique par les notables, n'aura en réalité profité qu'à quelques familles, mettant en lumière les tensions cachées de la société beaufortaine. »
	\footcite[32]{HViallet1993}
\end{quote}


\section{Les ventes des communaux}

\begin{quote}
	« Les acensements des années 1760
	
	L'acensement d'une partie des communaux, qui constitue la première étape d'une volonté de démembrement de la propriété collective, met en lumière, derrière les troubles suscités par les opposants, les luttes et les ambitions des différents acteurs de la communauté beaufortaine. Le conseil de Saint-Maxime en prend la décision le 11 février 1761. « Envisageant le peu de profit qui revient à la communauté de la grande étendue de communaux qu'elle possède et dont ne jouissent que les particuliers aboutissants... qui d'ailleurs n'en retirent pas, à beaucoup près, l'avantage qu'ils en auroient s'ils étoient économiquement distribués et un certain superflu d'iceux retranchés ; sans endommager personne, l'on pourroit par ce moyen des acensements faire un revenu à la communauté » (A.C.B., BB2, Registre de délibérations du conseil 1754-1771, p. 70).
	
	Le 12 mars, le conseil obtient l'autorisation de l'intendant, conformément à l'article 57 de la Péréquation Générale (A.C.B., DD 4, Délibération du conseil du 28 mai 1762). Les pâturages destinés à la location sont choisis (A.C.B., DD 194, Ventes de communaux 1768-78. Délibérations, ordonnances, verbaux), il s'agit de terrains isolés et éloignés, d'altitude élevée :
	
	— 300 jx à Outray, « vallée d'une montée rapide et scabreuse, et la moins à portée du public, n'y ayant que deux particuliers qui y possèdent des fonds. Les autres habitations et hameaux sont éloignées de plus deux heures de chemin » ;
	
	— l'Arpette au sommet de la vallée de la Gittaz (environ 70 jx) : « le bas de la vallée est éloigné du sommet de trois heures, et de cinq heures du chef-lieu. (Les communs) ne sont praticables par aucun hameau, ils ne servent qu'à l'usage des particuliers qui y ont des montagnes, qui sont huit dont deux n'en profitent pas pour n'être pas à leur portée » ;
	
	— 80 jx à Plan Vorex entre la Gittaz et Roselend. « Cette commune n'est à la portée que de trois familles » ;
	
	— 37 jx à la Grande Lanche de la Gittaz ;
	
	— 63 jx à la Luaix sous le Roc du Vent ;
	
	— 100 jx à Pra Lochet (vers la Berge) ;
	
	— 40 jx au Crêt du Gex (au-dessus de Roselend) ;
	
	— 20 jx aux Collombières (au-dessus d'Arêches) ;
	
	— plus de 100 jx au Mirantin ;
	
	— 54 jx au Glaçon (vallée de Treicol).
	
	Un mois plus tard, s'élèvent les premières oppositions, sous forme d'une requête adressée à l'intendant (A.C.B., BB 2, p. 78, Requête du 9 avril 1761), qui suspend la procédure des acensements. Après une réplique du conseil et de nouvelles demandes, une ordonnance de l'intendant du 19 mai 1762 (A.C.B., BB 2, p. 96, Ordonnance du 19 mai 1762) permet la reprise des opérations. Aussitôt après, les opposants déposent une nouvelle requête (A.C.B., DD 4, Requête du 28 mai 1762), mettant ainsi un sursis à l'ordonnance. Les enchères sont suspendues. Au cours de l'été 1762, les opposants envoient à dessein leurs troupeaux sur les communaux destinés à l'acensement, d'où plusieurs procédures entamées contre eux (A.C.B., DD 4). Au début de 1763, les terrains à louer sont définitivement fixés ; les enchères ont lieu d'avril à mai, après inhibitions faites par l'intendant de les troubler (A.C.B., DD 7, Propriété communale. Police rurale). »
	\footcite[70--71]{HViallet1993}
\end{quote}

\begin{quote}
	« Ce sont les notables qui vont profiter des ventes en agrandissant leurs alpages : 
	\begin{quote}« les six plus gros achats, réalisés par huit personnes, recouvrent à eux seuls 60 \% des surfaces aliénées. Quelques autres acheteurs agrandissent leur domaine privé de 20 à 50 ha et s'approprient 16 \% de la superficie vendue. Par contre les petits acheteurs, les plus nombreux n'acquièrent qu'une partie minime des biens communaux. 85 achats, dont 59 individuels concernent des surfaces inférieures à 5 ha et ne représentent que 12 \% des aliénations. Cette très forte concentration des achats entre quelques mains est à rapprocher de l'importante différenciation sociale existant à Saint-Maxime »»%
		\footnote{Cité par \footcite[79]{HViallet1993}, renvoyant à \footcite[188]{Fleury1985CommunalBeaufort}.}
	\end{quote}
	. [...] le secrétaire Joseph Blanc et Me Jean-Baptiste Blanc achètent à eux seuls 25 \% des surfaces aliénées. »
	\footcite[79]{HViallet1993}
\end{quote}


\begin{table}[htbp]
	\centering
	\caption{Les acquéreurs de biens communaux}
	\label{tab:acquereurs_communaux}
	\begin{tabular}{lrrr}
		\toprule
		Catégories sociales & Surface achetée & \% surface achetée & Zone d’alpages (\%) \\
		\midrule
		Me Joseph Blanc & 99 ha & 14,3 & 17,8 \\
		Me J.-B Blanc & 75 ha & 10,9 & 14,8 \\
		Propriétaires exploitants d’alpage (7) & 282 ha & 40,9 & 26,2 \\
		Fermiers d’alpage (2) & 5 ha & 0,7 & 0 \\
		Professions non agricoles aisées & 103 ha & 14,9 & 16,7 \\
		Paysans du bas du finage & 102 ha & 14,8 & 24,5 \\
		Pauvres & 24 ha & 3,5 & 0 \\
		\midrule
		Total & 690 ha & 100 & 100 \\
		\bottomrule
	\end{tabular}
	
	\footnotesize
\end{table}
\emph{Source :} Hélène VIALLET\footcite[80]{HViallet1993}, d’après Philippe FLEURY \footcite[40]{Fleury1985CommunalBeaufort}


\begin{quote}
	« Le conseil manifeste une volonté ferme de poursuivre, sinon les ventes, du moins les acensements. Dès le 11 juin 1773, il demande une autorisation d'acenser les communaux invendus : « cet acensement servira non seulement à produire un revenu à la communauté en extinction des intérêts dûs pour les affranchissements, mais encore à insinuer au peuple que l'on ne se désiste point du projet de la vente de ces parcelles, et à le désabuser de l'idée qu'il a pu concevoir, qu'au moyen d'une ligue formée pour ne point les miser, elles resteroient communes, et qu'il en jouiroit comme par cy-devant »(A.C.B., DD4, Délibération du conseil du 11 juin 1773). Soit par cabale, soit par manque réel d'intérêt, ces terrains ne sont pas vendus en 1776, lors de la deuxième opération d'enchères. Le 5 avril 1777, le conseil préfère renoncer à la vente de ces treize articles, pour « éviter autant que possible au public tout sujet de plaintes sur la privation des communes à laquelle il paroît si sensible »(A.C.B., BB 3, Registre de délibérations du conseil, 1771-1784, p. 85). Une seule de ces parcelles est de taille importante, il s'agit de Rognoux ; cette montagne servira « pour tenir durant l'été au paquéage les troupeaux qui pourroient se trouver infectés de maladie communicative ». »
	\footcite[81]{HViallet1993}
\end{quote}

\section{Blanc encore}




\chapter*{Conclusion}
\addcontentsline{toc}{chapter}{Conclusion}

	Naissance d’un capitalisme agricole
	Limites et questions en suspens
	-	Propriétés dans le mandement et au-delà
	-	Les stratégies matrimoniales et morcellement
	-	Pourquoi continuer à faire venir des suisses ?
	
	
	
	Texte avec une personne \idxnom{Cornu, André}
	puis une autre personne \idxnom{Blanc, Michel}
	et un lieu \idxlieu{Beaufort (Savoie)} ou \idxlieu{Hauteluce}.
	

	
	\section{Un tableau}
	
	\begin{table}[htbp]
		\centering
		\begin{tabular}{llr}
			\toprule
			Mas & Nature & Surface (ha) \\
			\midrule
			X & Herbage & 12.3 \\
			Y & Bois    & 4.8  \\
			\bottomrule
		\end{tabular}
		\tabcaptionindex{tab:surfaces}{Surfaces par nature}
	\end{table}
	

	
	\appendix
	\chapter{Index}
	
	\printindex[nom]
	\printindex[lieux]
	\printindex[auteur]
	\printindex[tab]
	\printindex[carte]
	\printindex[image]	
	\nocite{*}
	\printbibliography

	\appendix
	
\chapter{Méthodologie}

\maketitle

\section{Introduction}


\subsection{problèmes du nombre de parcelles}

Comme nous l’avons vu, le fichier Excel Barbéro contient 17803 parcelles.

Le numéro maximal attribué à la dernière parcelle répertoriée dans les documents de la Mappe de Beaufort aux ADS est : 17691.

Enfin le nombre de numéros de parcelles géoréférencées extraites des cartes Barbéro est seulement de 11787.

Une étape préalable consiste donc à comprendre ces différences.

a)	Nous avons constaté dans le fichier Beaufort-FICHIER-ATLAS-OK.xls qu’il manquait 136 numéros dans la séquence continue de 1 à 17691 ; en consultant les registres aux ADS, et en effectuant un sondage, il s’avère que ces numéros « manquants » y sont rayés et il est fait mention que ces parcelles ont été transférées  à la commune de « La Chapelle » à un moment donné durant l’élaboration de la Mappe ; « La vallée de Treicol est un secteur âprement disputé par Aime et les Chapelles, contre Saint-Maxime. Elle fait l’objet d’un procès pluriséculaire, puisqu’il débute en 1599 par le refus de Me Jean-François Dunand de payer la taille de sa montagne de Presset à Beaufort. L’affaire s’envenime, car de nombreux particuliers possèdent des alpages à Treicol. Pour être limitrophe de la vallée et augmenter ses prétentions, la communauté des Chapelles va jusqu’à acheter en 1732 une longue et étroite bande de rochers depuis la Pierre-Menta jusqu’à son territoire. Portée devant toutes les instances possibles, l’affaire n’est définitivement tranchée qu’en 1907, par un arrêt du Conseil d’État qui attribue la vallée à Beaufort.» 	\footcite{HViallet1993}.

La Chapelle est devenue la commune actuelle de la cote d’Aime (fusionnée en 2016 dans la Plagne Tarentaise). 

Carte9 : la zone « inconnu » correspond à l’information du « mas » des parcelles, ces parcelles n’étant pas référencées dans le fichier Excel Barbéro. 

Nous avons ajouté au fichier les 136 parcelles manquantes, ces lignes ne comportent aucun attribut descriptif.

b)	Autre particularité du fichier Excel Barbéro est qu’il contient :

-	Une parcelle non numérotée, mais avec la lettre « A » dans le champ n° de parcelle.

-	6 parcelles différentes, mais dont les numéros sont en doublon.


Récapitulatif :

Analyse des numéros 
du tb excel Barbéro	Tableau
Barbéro	Nb Entiers
+ Parcelles
manquantes	Hypothèses
Parcelles avec nb entier entre 1 et 17691	17552	17552	
Parcelles avec le même n° (nombre entier)	2	1	le doublon est une décimale omise
Parcelles manquantes (Transférées)		136	
Parcelle numérotée "A'	1	0	le n° omis est un nombre décimal
Parcelles avec un n° + une lettre	2	2	
Parcelles avec n° + lettre (n° déjà existant)	10		
Parcelles avec une décimale (1/2 ; 1/3,…)	232		
Parcelles avec le même n° nb + décimale)	4		
Total	17803	17691	

Conclusion : si nous totalisons uniquement les numéros avec un nombre entier, nous retombons sur le dernier numéro de parcelle (17691) des registres de la Mappe de Beaufort.


\subsection{Import de la Mappe dans QGIS}

La première étape \footnote{projet QGIS : Assemblage précis.qgz} a consisté à superposer à plat les différents éléments disponibles Cette approche évite à ce stade le  risque de déformation du fait d'un géoréférencement d'une carte ancienne sur une carte moderne IGN) : 

- calage les unes par rapport aux autres, les quatre images de la vue globale de la Mappe Sarde

- superposition par-dessus (avec l'outil Georeférencer de QGIS) les 45 photos détaillées de la Mappe téléchargées des ADS. 



\cartefigure
{figures/1 MS 45 photos détails non GeoRef.jpg}
{fig:MS45nonGeoRef}
{45 photos assemblées de la Mappe Sarde de Beaufort}

\cartefigure
{figures/1 Assemblage des 45 photos détail Mappe Sarde Beaufort.jpg}
{fig:MS45Assemblage}
{assemblage des 45 photos de la Mappe Sarde de Beaufort}

Remarques :\\
- une photo est manquante sur le site des ADS (\texttt{361\_b\_16}). Afin de la constituer, nous avons découpé un morceau de  la photo globale associée . Bien évidemment la qualité est beaucoup plus faible.

- Le collage des 45 photos a fait apparaître quelques non-recouvrements du fait du photographe des ADS

- la Mappe Sarde de Beaufort est parfois très abîmée en particulier la jonction entre les deux toiles peintes est impossible du fait des manques.

\cartefigure
{figures/1 MS 45 photos détails Anomalies.jpg}
{fig:MS45_Anomalies}
{photos 361_013 et 014 non jointives avec 019 et 020}






\subsection{Retraitement des cartes Barbéro}

Le CD-Rom fournit 55 images-cartes comme celle-ci :

\cartefigure[0.5\textwidth]
{figures/1 exemple carte CD Barbéro.JPG}
{fig:BarberoCarte01}
{exemple Carte 01 CD Barbéro} 

Telle quel, ces cartes-images, ne sont pas plus exploitable que les photos de la Mappe déjà importées.
En revanche, elles contiennent deux éléments seulement : des lignes et des chiffres.

Des outils permettent des retraitements suffisamment automatiques pour rendre possible une alternative au travail manuel sur les 55 images et les 17691 parcelles. Le détail des outils et des étapes est décrit en annexe (cf Annexe3), en voici les tâches principales :

-	Détourage des 55 photos pour ne garder que les lignes et les numéros

-	Séparation des figures : 55 images avec seulement les contours des parcelles et 55 images avec les numéros des parcelles seulement .

\cartefigure[0.5\textwidth]
{figures/1 Beaufort secteur 01 original.jpg}
{fig:BarberoCarteNumOnly}
{Détourage et conservation des numéros exemple Carte 01}


\cartefigure[0.5\textwidth]
{figures/1 Beaufort secteur 01 reds.jpg}
{fig:BarberoCarteRed}
{Détourage et conservation des contours des parcelles exemple Carte 01}



-	Décodage des chiffres, les outils d’OCR que nous avons essayés ne donnant pas de résultat, nous avons dû coder en python une reconnaissance de formes au niveau pixel ; chaque un des nombres (0 à 9) ayant plus ou moins une signature en pixel unique.

-	Nous avons gardé la position des ensembles de pixels composant chaque chiffre. Puis nous avons regroupé les chiffres mitoyens pour aboutir a) à un nombre = au numéro de la parcelle et b) une référence (X,Y) du pixel au centre du groupe de chiffre composant un nombre.

-	Un retraitement manuel sur une centaine de nombres non reconnus a été nécessaire : nombre tronqué en bordure d’image, nombre mélangé à des éléments superposés comme la boussole ou la barre d’échelle ou encore chiffre mal pixélisé non reconnu dans le traitement précédent.

Au final de cette étape nous obtenons :

-	55 images des contours de parcelles 

- 	55 images des numéros (images inutilisées sauf pour vérification)

-	55 tables avec les nombres reconnus et leur position pixel (X,Y) sur chaque image.

Les 3 ensembles sont exactement superposables

Nous avons rassemblé les 55 tables en un seul fichier \texttt{XY\_No\_Parc.csv} (fichier ; X ; Y ; No\_Parc) et nous avons éliminé les doublons (parcelle se trouvant affichée sur deux images de secteur mitoyen). Au final nous avons 11.744 numéros distincts reconnus. Ce nombre représente les 2/3 (66\%) des 17.803 parcelles référencées dans la table Barbéro.

C'est assez compréhensible : la table Barbéro (fichier Excel) est issue des registres associés à la Mappe alors que les numéros sur les cartes sont issus de la lecture de la carte elle-même où les numéros ne sont pas toujours lisibles.

Nous verrons comment localiser approximativement les parcelles du tiers restant avec une méthode fondée sur la localisation du Mas de chaque parcelle. 




\section{Attribution d’une altitude à chaque parcelle}

\subsection{Géoréférencement des cartes Barbéro}

Le CD-Rom contient une carte d’assemblage des 55 secteurs :


\cartefigure 
{figures/1 Tableau d'assemblage des secteurs.JPG}
{fig:BarberoAssemblage}
{Carte d'assemblages de cartes fournie dans le CD Barbéro}


En utilisant la fonction géoréférencement de QGIS, nous avons superposé les 55 images des contours des parcelles sur l’assemblage détaillé de la Mappe. Cette opération était relativement simple (mais longue !) puisque nous pouvons assez facilement faire correspondre des points de ce géoréférencement en retrouvant les lignes des parcelles des cartes Barbéro avec celles correspondantes de la Mappe.

Malheureusement dans un ou deux cas l’assemblage des cartes Barbéro comportait des lacunes (certaines cartes ne se superposaient pas). C’est dommage, mais l’impact est resté très marginal.

Ensuite, nous avons réutilisé les mêmes points de géoréférencement de l’étape précédente, pour géoréférencer les cartes des numéros de parcelles, afin que chaque numéro se retrouve au bon endroit.

À ce stade, nous avons 3 couches superposées avec :\\
- les photos détaillées de la Mappe\\
- les contours des parcelles\\
- les numéros (en mode image, inexploitables mais utile pour vérifications ultérieures)


La dernière étape a consisté à transformer la table des numéros de parcelles en une couche vectorielle avec une table d’attributs obtenant ainsi la géolocalisation des numéros de parcelles.

un programme python qui : \\
- prend chaque la table "i"\\
- la charge en layer vecteur (add delimited text)\\
- l'exporte en fichier gpkg (format de couche vectorielle de QGIS)\\
-charge ce fichier dans le GeoRefrencer de QGIS et l'associe au fichier des points utilisé pour géolocaliser les contour des parcelles\\
- lance le géoréférencement

Au final nous avons 55 fichiers gpkg qui se superposent aux trois couches raster précédentes.

Pour chacune des couches je fusionne tous les éléments en un seul :\\
- les 45 images raster de la Mappe en un raster virtuel : \texttt{Virtual photo detaillées.vrt} \\
- les 55 images des contours des parcelles en un raster virtuel \texttt{Virtual_Red.vrt}\\
- les 55 fichiers gpkg des numéros de parcelles en un seul gpqk  \texttt{text}\\

Remarque : il y a dans le fichier des numéros de parcelles des numéros en double (ou triple et même quadruple. Ce sont sans doute des numéros qui apparaissent sur plusieurs secteurs Barbero (les secteurs se chevauchent) mais il y a peut-être d'autres anomalies. Nous avons éliminé tous les doublons à moins de 100m les uns des autres. Il reste 35 anomalies de doublons éloignés l'un de l'autre. Il s'agit sans doute de reconnaissance erronée des numéros sur la photo de la Mappe par D. Barbero (exemple sur le secteur 44 une parcelle est lisible sur la carte comme la numéro 16875 mais elle est transcrit par sur par D.Barbero en 6871), l'autre explication possible une reconnaissance erronée des numéros sur les images Barbero par notre programme de retranscription.




Le décodage numéros de parcelles à partir des cartes de secteur Barbéro a abouti à reconnaître 11787 numéros de parcelles (et à les géoréférencer).

Les deux tables, celle du fichier excel du CD-Rom (enrichie par notre modélisation) et celle issue du géoréférencement des numéros de parcelles ont un index commun qui est le numéro de parcelle (dans quelques cas, une dizaine, une lettre est associée au numéro d'une parcelle, dans ce cas une décimale a été ajoutée au numéro de la parcelle par exemple la parcelle 17301 A a été numérotée 17301.02 cette décimale permet de ne pas confondre avec les décimales de 0.5 pour 1/2 et 0.33 pour 1/3 ). 

En fusionnant les deux fichiers, il est donc aisé de créer une seule carte vectorielle avec des données attributaires pour toutes les parcelles et une géolocalisation pour les 2/3 d'entre elles.

Mais il manque 6016 parcelles non géoréférencées, ce qui représente tout de même un tiers des parcelles.

Nous ne savons quels ont été les moyens et méthodes de Dominique Barbéro pour faire son travail de cartographie: Est-ce parce qu’il n’a pas pu lire certains numéros sur la Mappe originale ? Est-ce seulement l’édition des images de secteurs qui n’a pas affiché tous les numéros et toutes les parcelles ?

Toujours est-il que la proportion de parcelles non référencées pose un vrai problème pour l’objectif de cette étude qui est d’isoler les parcelles d’altitude.

\subsection{Geolocalisation précise sur carte Ign}

Nos trois couches :  photo de la Mappe (Raster), contour parcelle (Raster) et Parcelles Barbéro (vecteur), se superposent parfaitement.

Nous avons alors géoréférencé précisément la Mappe,  sur un fond de carte IGN en utilisant une variété de points de calage : \\
- les contours de la commune (en particulier les angles )\\
- les églises, chapelles, village...\\
- les cours d'eau\\
puis nous avons réutilisé ces points de calage pour géoréférencer les contours de parcelles et surtout la couche des parcelles.

\cartefigure
{figures/1 Mappe de Beaufort Georéférencée sur IGN.jpg}
{fig:MSsurIGN}
{Mappe de Beaufort Georéférencée sur IGN}

Le résultat est tout à fait satisfaisant, dans le sens où il n'y a pas de déformation exagérée de la Mappe de Beaufort quand on la transforme pour qu'elle se cale sur la carte moderne IGN

\subsection{Geolocalisation des parcelles manquantes}

Il reste, à ce stade 6016 parcelles non géoréférencées, mais nous connaissons les mas où elles se situent.

Les mas sont des territoires à l’intérieur du territoire de Beaufort, pour la plupart il s’agit d’un hameau ou groupement de maisons et les parcelles qui en dépendent, autrement ce sont des lieux dits. Toujours est-il que toutes les parcelles (le 17803) sont associées avec un « mas » et un seul.

Le CD-Rom Barbéro contenait une carte des mas, mais nous avons préféré reconstituer nous-mêmes les mas avec l'outil cartographique QGIS.

Nous avons donc créé une couche vectorielle des Mas, en créant des polygones qui entourent toutes les parcelles déjà géoréférencées.

Ce faisant nous avons constaté un certain nombre d'anomalies, en particulier des ensembles disjoints d'un même mas.

Nous avons alors créé un attribut "Mas2" qui est égale à "Mas" sauf dans certain cas, par exemple nous avons crée "Acrays 2" , "Ban 2", "Curtillet ?", "Dommellin ?", "Isle et Lacura ?", "Pecherie", "Plan ou Plains ?", "Praz" et "Territoire transféré à LaChapelle". Le principe était de n'avoir que pour chaque mas (Mas2) un territoire homogène.

\cartefigure
{figures/1 Carte des Mas.jpg}
{fig:Mas2}
{Carte des Mas reconstitués}

xxxxxxxxxxxxxxxx cas des parcelles hors de leur mas xxxxxxxxxxxx



\subsection{Ajout de l'altitude des parcelles}

Les cartes IGN fournissent aussi des couches de courbes de niveau, permettant de déterminer l'altitude de chaque point. Nous avons donc ajouté deux informations à la table attributaire des parcelles : l' altitude et une tranche d'altitude : \texttt{ [inf à 800m], [800-1000m], [1000-1200m], [1200-1400m], [1400-1600m], [1600-1800m] et [sup à 1800m]}.

Nous avons aussi ajouté une couche matérialisant les zones où l'altitude est supérieure à 1600m.

Il n’y a pas de définition précise des zones d’alpages par rapport aux zones d’hivernage permanent ou aux zones de pré-alpages. Dans le système agropastoral dit de « grande montagne » \footcite{ArbosMouthon2017} \footcite{Arbos1922}, il s’agit des pâturages utilisés entre la Saint XXX et le 24  juin (fête de la Saint Jean-Baptiste) et le 12 septembre (fête de la Saint XXX). 

Nous savons qu’à la fin de l’été, les grands troupeaux sont démantelés et dispersés dans des chalets capables de loger entre 5 et 10 vaches. Les vaches des grands troupeaux d’alpage pouvant appartenir soit au propriétaire (ou tenancier) de l’alpage soit au propriétaire (« client » de l’alpagiste) du chalet d’hivernage.

Suivant la richesse (en terre) des paysans, ils pouvaient posséder 3, 4 … 6 chalets intermédiaires entre la zone estivale et la zone d’hivernage finale en bas de vallée.

L’architecture des chalets varie en fonction de leur utilisation. Il est facile en se promenant dans la vallée du Beaufortain d’identifier à quelle altitude, les chalets d’hivernage disparaissent et sont remplacés par des granges, sans étable ni fenêtre ni espace d’habitation permanente. Par ailleurs une partie non négligeable des vieux chalets datent du 18e siècle. 

Cette approche empirique nous a conduits à fixer à 1600m l’altitude au-delà de laquelle on peut parler d’alpage. Nous avons identifié les derniers chalets d’habitat permanent à 1500m mais nous avons supposé que les prés ou pâturages avoisinant ces chalets étaient utilisés pour récolter le foin pour passer l’hiver et non pour la grande pâture.

Cette hypothèse est confirmée par l’exploitation actuelle des alpages, bien qu’il faudrait aussi prendre en compte l’histoire du climat. Le XVIIIe siècle était nettement plus froid qu’au XXe siècle où s’est mis en place l’agropastoralisme sous sa forme actuelle. Par conséquent la marge de 100m par rapport aux dernières habitations permanentes nous paraît aussi légitime de ce point de vue.




\subsection{Point d'étape : toutes les parcelles sont situées géographiquement}

Jusqu'à présent, le travail fut très technique, long et parfois fastidieux, mais cela nous a permis d'obtenir enfin une base de travail utilisable en Histoire.

Remarque : ce travail est déjà fait et disponible en ligne pour les communes de Hautes Savoie par les Archives Départementales de Haute-Savoie. Pour la Savoie, il est fait partiellement (CD Barbéro) et sur un nombre de communes limité. Sur la vallée du Beaufortain, seule la commune de Beaufort a été traitée au format Barbéro, rien n'est disponible pour les communes de Queige, Villard-sur-Doron et Hauteluce en dehors des documents d'origine de la Mappe. 


A ce stade nous avons donc une nouvelle table reprenant toutes les informations de la table initiale du CD Barbéro avec en plus pour chaque parcelle :\\
- son géoréférencement \\
- son altitude et sa classe d'altitude\\
- son occupation du sol plus utilisable que l'attribut nature du sol\\
- sa surface en ha (déduite de celle en m²)\\
- son mas (mas2) retravaillé pour n'avoir que des territoires d'un seul bloc \\
- des informations plus techniques comme le nom du secteur Barbéro sur lequel se situe la parcelle et ses coordonnées en pixel, et pour les parcelles n'en faisant pas partie la méthode qui a permis de les placer par rapport aux autres)



\section{Constitution d'une base de données enrichie}

Afin de rendre possible l'exploitation historique du cadastre, les données initiales ont été complétées par plusieurs informations nouvelles.

Les graphies manifestement variantes d'un même mas ont été harmonisées afin de constituer des ensembles territoriaux cohérents. Un attribut \texttt{Mas2} a ainsi été créé pour faciliter les traitements statistiques et cartographiques.

Les nombreuses dénominations cadastrales de la rubrique « Nature » ont également été regroupées dans une nomenclature simplifiée d'occupation du sol. Cette opération avait pour objectif de conserver la richesse descriptive des registres tout en permettant une lecture synthétique de l'organisation du territoire.

Enfin, chaque parcelle a été associée à une altitude obtenue par croisement avec un modèle numérique de terrain. Des classes altitudinales ont ensuite été définies afin de permettre l'étude de la répartition des propriétés, des usages du sol et des formes d'exploitation selon l'altitude.

\section{Résultat obtenu}

À l'issue de ce travail, chaque parcelle dispose :

\begin{itemize}
	\item des informations cadastrales issues des registres sardes ;
	\item d'une localisation géographique ;
	\item d'une appartenance à un mas ;
	\item d'une catégorie homogène d'occupation du sol ;
	\item d'une altitude et d'une classe altitudinale ;
	\item d'une superficie exprimée en hectares.
\end{itemize}

Cette base de données constitue le corpus utilisé dans l'ensemble des analyses présentées dans les chapitres suivants.

Les traitements techniques ayant permis la constitution de ce corpus sont présentés en détail dans l'annexe méthodologique.

\section{Méthodologie}

\subsection{Analyse de la table du CD-Rom : numérotation des parcelles}

Le CD-Rom Barbéro contient un ficher Excel : Beaufort-FICHIER-ATLAS-OK.xls contenant 17803 parcelles  (avec en plus, une ligne d’en-têtes des colonnes) :

•	Numéro : il s’agit des numéros de parcelles a priori numériques dont le plus grand numéro est 17691. Il devrait a priori avoir 17691 parcelles sur la commune de Beaufort du Cadastre Sarde mais on constate immédiatement que le nombre de lignes ne correspond pas au nombre de parcelles

•	CF : Chef-lieu « O » ou « null » ; signification à investiguer ultérieurement

•	E : Ancien fond ecclésiastique « O » ou « null » ; signification à investiguer ultérieurement

•	I : « O » ou « null » ; signification à investiguer ultérieurement (dans la notice livrée avec le CD-Rom il est fait mention à cet emplacement d’une colonne nommée « N » (pour Noble)

•	Propriétaires : nous avons là des noms et prénoms des propriétaires avec parfois des compléments d’information par exemple « MONOD Jeanne soit Popelin veuve d'Antoine GACHET » ou « TARIN Nicolas et Humbert son frère» et aussi des noms de communautés par exemple « La communauté pour l'usage commun » ou « Chapelle de Roselend »

•	Statut : cette colonne contient du texte avec des valeurs répété : « Bourgeois » « communier » … il s’agit d’une information sur le statut du propriétaire a priori tel qu’il est inscrit sur le registre associé au cadastre

•	IND : parcelle en indivision, elle contient des valeurs « null » et quelque « O »

•	Mas : 

•	Nature : cette colonne contient du texte avec des valeurs récurrentes : « Roc », « Pâturage » ; « Marais et pâturage » … qui sont des descriptions de la nature du sol

•	DB : cette colonne contient des valeurs numériques 0, 1 , 2 ou 3 qui correspond aux quatre « degrés de bonté ». Le degré de bonté est une évaluation de la qualité (ou du rendement) du sol, faite par les agents du duché de Savoie en charge du relevé du cadastre, il a un but fiscal

•	m² : il s’agit sans doute de la surface de la parcelle en m² ; le registre du cadastre contenant de valeur de superficie selon une double mesure : la mesure de Savoie (Journaux, toises, pieds) et mesure de Piémont (Journaux, Tables, Pieds). 1 journal de Savoie équivaut à 0,3685 ha, celui du Piémont lui est légèrement supérieur. Cette conversion a été réalisée sans doute par le créateur du fichier D.BARBERO. (remarque : l’ensemble de la commune actuelle couvre 149,53 klm²  mais la partie transférée à l’époque de l’établissement de la Mappe Sarde est revenue à la commune actuelle de Beaufort sur Doron, la somme des surfaces du fichier D.BARBERO est égale à 139.648.325 m² soit 139,6 klm² ce qui est cohérent, on peut donc valider que la surface sur fichier Excel est en m² d'autant plus qu'à l'époque de l'établissement de la Mappe une partie des parcelles a été rayée car réattribuée à la paroisse de LaChapelle, elles ont été depuis réintégrées à la commune de Beaufort, dont il est normal que la superficie calculée par Barbéro soit inférieure à la superficie de la commune actuelle de Beaufort-sur-Doron) 

•	Avoine : des valeurs null ou numérique s’étalant de 13,5 à 22

•	Seigle : des valeurs null ou numérique s’étalant de 0,5 à 4 (les valeurs 0,5 sont stockées au format texte « 1 / 2 » nous les avons changées en 0,5

•	F. Bœuf : sous bœuf est regroupé le nombre de bovins, la colonne contient des valeurs numériques de 0,5 à 9 puis pour les valeurs supérieures le nombre est suivi de la lettre « L » : « 10 L », « 12 L », « 100 L » …

•	F. Cheval : sous Cheval est regroupé le nombre d’équidés de toute sorte (âne, mule, mulet cheval ) . La colonne contient des valeurs numériques de 0,5 à 9 puis pour les valeurs supérieures le nombre est suivi de la lettre « L » : « 10 L », « 12 L », « 100 L » …

•	Fascines : (fagot de petit bois ?) valeur numérique non investiguée pour l’instant

•	Chevron : valeur numérique sans doute bois d'œuvre pour charpente

•	Latte : valeur numérique sans doute bois d'œuvre pour des planches

•	Pièce 1 : valeur numérique non investiguée pour l’instant

•	Pièce 2 : valeur numérique non investiguée pour l’instant

•	N° Grief : champs pouvant contenir une ou plusieurs valeurs numériques, il s’agit de numéro de référence des contestations faites auprès des autorités relativement à la parcelle cadastrée. En principe il faudrait normaliser ce champ et créer une logique de relation 1-n entre une parcelle et les griefs associés, cependant il n’y a que 22 cas où une parcelle est associée à 2 ou 3 griefs, et l’étude des griefs est très secondaire dans notre étude.

•	Grief : champ texte contenant une typologie selon les types de griefs

•	Problèmes : descriptif sommaire du grief


Degré de Bonté :

Dans l’attribution des degrés de bonté, chaque type d’occupation du sol comporte les trois classes ; la qualité de la parcelle pouvant être bonne (1), moyenne (2) ou mauvaise (3). Par exemple pour la paroisse de Landry en vallée de Tarentaise (Savoie) la « note des estimes en degré avec leurs catégories… » nous informe sur la production qu’un terrain doit fournir pour se voir attribuer tel ou tel degré de bonté. Un pâturage doit produire annuellement deux quintaux de foin pour obtenir le degré de bonté 1, un quintal et demi pour le 2 et un quintal pour le 3. Pour les champs, la mesure est en bichet et il en faut vingt-neuf répartis en froment, seigle et farine pour répondre aux critères du degré de bonté 1, vingt-six pour le 2 et vingt-trois pour le 3, et ainsi de suite pour chaque usage du sol. Pour le bâti, les estimateurs attribuent le degré de bonté d’une terre voisine, de préférence appartenant au même propriétaire. Cette méthode d’évaluation des terrains fixe des niveaux d’exigence précis basés sur le rendement de la parcelle. Ainsi, les parcelles classées en 1 sont de bonnes terres productives, car les critères exigés pour avoir ce degré de bonté sont du même niveau. Le fait que chaque nature de sol comporte l’intégralité des classes (1, 2 et 3) et que les critères d’évaluation répondent à un même niveau d’exigence permet une comparaison entre deux terrains du même degré de bonté. \footcite{Baud2010} 


Mesure de Savoie :
« La base de ce système est le « pied de chambre de 0,033m qui élevé au carré donne le « pied de cadastre de 0.92m². un carré de 8 pieds forme la toise carrée (7.370m²) et 400 toises carrées équivalent à un journal de Savoie dont la valeur exacte est de 29.4838 ares. On compte pour les approximations 3 journaux à l’hectare » 	\footcite{GUICHONNET1955}


Mesure de Piemont :

« Les cadastreurs piémontais se servirent des unités de leur pays et qui dérivaient du pied liprand de 0,513 m. Il fallait 6 pieds pour faire un trabuc (3,082 m) et 2 trabucs pour une perche (6,164 m). Ces mesures linéaires servaient de base aux unités agraires suivantes : le pied de table (rectangle de 1 x12 pieds liprands de côté) ou 3,167 m². La table, ou perche carrée vaut 12 pieds de table (38,009 m) et 100 tables, ou 400 trabucs carrés, donnent un journal de Piémont. La mesure de Piémont est à préférer lorsqu’on utilise l’ancien cadastre. Pour les usagers soucieux de l’exactitude des contenances, elle est évidemment meilleure car elle résulte des opérations sur le terrain ; pour l’historien, elle est commode car la conversion des tables en journaux repose sur la base 100. » Idem


Mas :

«On a la preuve, par les reconnais; féodales, que les mas sont antérieurs de trois ou quatre siècles au cadastre, et que les « trabucants » piémontais n’ont fait qu’enregistrer un état de fait profondément et obscurément enraciné dans la tradition paysanne. Les mas appartiennent à plusieurs propriétaires, nobles ou roturiers, et le régime foncier de 1730 recouvre une réalité agraire très antérieure; ils chevauchent parfois plusieurs communes et ne portent pas toujours le nom des villages ou des lieux-dits qu’ils renferment, mais un toponyme d’étymologie romaine ou germanique. Leur étendue comprend rarement une seule zone de végétation, mais englobe différents types de terroirs et semble bien représenter une unité élémentaire d'exploitation, domaine d’un seul individu. » Idem p296

a) Nettoyage des données

•	Pb colonne « Numero »

En fin de fichier il y a 13 lignes avec des valeurs alphanumériques : 12 parcelles ont une lettre à la suite du numéro de parcelle et 1 parcelle a une lettre sans numéro :

17160 B;;;;BOUCHAGE Jean Baptiste;..\\
17161 A;;;;BUGAND Claude;… \\
17301 A;;;;MARTIN Claude;… \\
17302 B;;;;MARTIN Claude;… \\
17303 C;;;;La communauté pour l'usage commun;Communaux;… \\
17304 D;;;;La communauté pour l'usage Commun ;Communaux ;… \\
17305 E;;;;La communauté pour l'usage commun;Communaux;…\\
17306 F;;;;La communauté pour l'usage commun;Communaux;… \\
17307 G;;;;La communauté pour l'usage commun;… \\
17308 H;;;;La communauté pour l'usage commun;…\\
17309 I;;;;La communauté pour l'usage commun;\\
17310 K;;;;La communauté pour l'usage commun;\\
A;;;;TOCCAT Daniel;Communier;;;;3;1543;;;;;;;;;;;;\\

Nous avons créé une nouvelle colonne : Complément_numéro.

Les lettres des 13 lignes sont déplacées dans cette nouvelle colonne, et la ligne sans numéro se voit attribuer un numéro « 0 ».

Il y a aussi des parcelles dont le numéro est suivi de 1/2 , 1/3 ou 1/4 , ces valeurs ont été remplacées par des décimales (0.5, 0.3, 0.25) ajoutées au numéro de parcelle.

Il reste à comprendre pourquoi le nombre de lignes est significativement différent de la valeur maximum de numéro de parcelle. Mais nous allons d’abord importer ce fichier dans une base de données de travail.

•	Pb colonnes « CF » , « E », « I » , « IND »

Les valeurs « O » dans les colonnes « CF » , « E », « I » , « IND » ont été remplacées par des « Oui » afin de permettre une caractérisation booléenne.

•	Pb colonne « m² »

147 ;;;;VIALLET Jacques fils de Benoit;Communier;;Praz;Chenevier;1;\#VALUE!;21;3;;;;;;;;;;\\
La valeur \#value est remplacé par la valeur « null » \\
•	Pb colonne « Nature »

La colonne Nature est longue et contient beaucoup de double description, par exemple « bois et pâturage » ou « prés et teppes », une option de normalisation aurait été de décomposer en une relation 1-N : une parcelle pouvant être reliée alors à deux natures « bois et « pâturage » ; mais n’étant pas certain de comprendre précisément l’intention du géomètre-cadastreur, nous avons pris l’option de laisser pour l’instant le champ nature tel quel, mais en revanche de créer un concept « Occupation_Sol » pour regrouper les nombreuses « Natures » en ensemble homogène en prenant pour principe lorsqu’une Nature est « multiple » de prendre la première comme la Nature principale ; par exemple « « bois et pâturage » » aura la valeur « Bois » pour « Occupation_Sol »

•	Pb colonne « Mas »

Pour les Mas il y a des mas dont les libellés se ressemblent, par exemple : « Coutafalliat » et » Coutassalliat » (les » s » et les «f » dans la graphie du XVIIIe), mais comme les numéros de parcelles ne sont pas forcément continus entre les deux orthographes, nous avons décidé de laisser les libellés de mas tel quel en attendant un travail futur de représentation cartographique des mas et parcelles.

•	Pb colonnes « Avoine », «Seigle », « Boeuf », « Cheval », « Fascines » , « Chevron » , « Latte» :

Les champs « Avoine », «Seigle », « Boeuf », « Cheval », « Fascines » , « Chevron » , « Latte» décrivent soit une quantité de production d’Avoine de seigle, de fagots de bois (Fascines ?) et c’est logique qu’elle soit associée à une certaine parcelle précise, mais pour les autres il s’agit d’un nombre de bêtes ; cette dernière valeur pose question, car cela devrait:

•	Être associé au propriétaire et correspondre à la taille de son troupeau (cela poserait alors d’autres questions, car pour la gestion des alpages (en été) les troupeaux sont constitués de bêtes appartenant en propre et des bêtes en location) ; cela ne semble pas être le cas, car pour un propriétaire possédant plusieurs parcelles le nombre de bêtes varie selon une logique non claire ; exemple :
Parcelles	Propriétaires	Mas	Nature	DB	Surface	Boeuf\\
7806	BLANC Jean Baptiste	Val de Treicol	Grange	2	67	180\\
7807	BLANC Jean Baptiste	Val de Treicol	Grange	2	181	180\\
7808	BLANC Jean Baptiste	Val de Treicol	Pâturage	2	31130	180\\
7774	BLANC Jean Baptiste	Val de Treicol	Grange	2	206	160\\
7776	BLANC Jean Baptiste	Val de Treicol	Pâturage et rocher	2	86210	160\\
8331	BLANC Jean Baptiste	Couecle	Pâturage et broussailles	2	65791	150\\
8323	BLANC Jean Baptiste	Couecle	Pâturage	2	147294	120\\
8324	BLANC Jean Baptiste	Couecle	Grange	2	79	120\\
7802	BLANC Jean Baptiste	Val de Treicol	Pâturage	3	10135	50\\
7803	BLANC Jean Baptiste	Val de Treicol	Pâturage	3	22085	30\\

•	Ou être associé à une ou plusieurs parcelles : il est évident qu’un troupeau qui « remue » change de parcelle régulièrement.
Nous avons donc conservé ces champs tels quels comme attribut de la parcelle en attendant de comprendre comment cette donnée fonctionne (il faudra sans aucun doute revenir aux registres originaux manuscrits associés au Cadastre Sarde conservé aux Archives Départementales de Savoie)

•	Pb colonnes « Avoine » et « Seigle »

Par exemple dans la ligne :

5299 ;;;;DUITTOZ GENIS;Communier;;Bersend;Champ;3;377;14 1/2; 1/2;;;;;;;;;;\\
Il est inscrit 14 1/2 pour 14,5 et 1/2 pour 0.5

En ouvrant le fichier avec Excel, les valeurs sont bien comprises par le logiciel comme des nombres décimaux, sauf lorsque la valeur est 1/2 qui reste comprise comme du texte

Les 668 cas des colonnes « Avoine » et « Seigle » ont été corrigés.

•	Pb colonnes Bœuf et Cheval

Il existe des « L » qui suivent les valeurs supérieures à 10 ; je ne sais pas à quoi cela correspond (sans doute pas l’unité de mesure Livre. Ces « L » sont supprimés afin d’avoir des colonnes numériques. (Il faudra tout de même vérifier ce qui est inscrit dans le registre.
Il y a aussi un cas où la cellule Bœuf contient une virgule ; elle est remplacée par « Null»

Pour la colonne Cheval il y a aussi 4 cas de « 1/2» remplacés par « 0.5»

•	Pb colonnes No Grief, Grief, Problème

Il existe des lignes dont le champ « Problèmes» n’a pas de numéro de Grief associé sans numéro., nous avons ajouté une numérotation allant de 800 à 952 et la description « Grief sans numéro ».

•	divers :
Nous avons aussi modifié certains titres de colonnes (remplacement des é , ², œ, espace…) et ajouté une colonne » Analyse » (vide) qui servira pour stocker les résultats des analyses des numéros de parcelles.

Désormais nous avons un fichier csv contenant des champs propres et homogènes :

Analyse	Texte\\
Numero	Numerique\\
Complement_numero	Texte\\
CF	booléen (Oui/Non)\\
E	booléen (Oui/Non)\\
I	booléen (Oui/Non)\\
Proprietaires	Texte\\
Statut	Texte\\
IND	booléen (Oui/Non)\\
Mas	Texte\\
Nature	Texte\\
DB	Numerique\\
Surface	Numerique\\
Avoine	Numerique\\
Seigle	Numerique\\
Bœuf	Numerique\\
Cheval	Numerique\\
Fascines	Numerique\\
Chevron	Numerique\\
Latte	Numerique\\
Pièce1	Numerique\\
Pièce2	Numerique\\
No_Grief	Texte\\
Grief	Texte\\
Problemes	Texte\\

\subsection{Constitution d’une table attributaire enrichie}

Un travail de modélisation de la table contenue dans le CD Barbero est disponible en annexe: ce travail permet de distinguer les entités : Mas, Propriétaire et Nature

Au niveau des Mas nous avons par exemple fusionné le mas "Coutassallat" dans le mas "Coutafaillat" ainsi que le mas "Treverse" dans le mas "Traverse"

Au niveau des nombreuses "natures", nous avons créé une entité "Occupation du Sol" afin de regrouper les 440 natures différentes.

Regroupement des natures cadastrales

Les registres de la Mappe sarde décrivent chaque parcelle au moyen d'une rubrique dite « nature ». Cette information est particulièrement riche et reflète la diversité des milieux exploités par les communautés montagnardes du Beaufortain. Outre les catégories classiques telles que les champs, prés ou bois, les estimateurs distinguent de nombreuses situations intermédiaires : pâturage sur roc, roc et pâturage, bois et broussailles, rocher broussailles, glières, rippes, marais, rivages, etc. Au total, plusieurs dizaines de libellés différents apparaissent dans les registres.

Cette richesse descriptive constitue un atout pour comprendre la perception et l'évaluation des terres par l'administration sarde. En revanche, elle rend difficile toute analyse statistique ou cartographique globale. La multiplication des catégories produit des tableaux peu lisibles et empêche de dégager les grandes structures de l'occupation du sol à l'échelle de la communauté.

Afin de permettre les traitements statistiques, les différentes natures ont donc été regroupées en un nombre limité de catégories fonctionnelles. Ce regroupement ne vise pas à remplacer les dénominations originales, qui demeurent conservées dans la base de données, mais à faciliter l'analyse des grandes composantes du territoire.

La méthode retenue repose sur une hiérarchie fondée sur l'usage économique principal de la parcelle. Lorsqu'un libellé comporte le terme « pâturage », la parcelle est classée dans la catégorie « Pâturage », quelle que soit la présence éventuelle d'autres éléments (roc, bois ou broussailles). Ce choix s'explique par l'importance économique de l'activité pastorale dans le Beaufortain et par le fait que de nombreuses parcelles d'alpage associent naturellement herbe, affleurements rocheux et végétation ligneuse. (remarque : nous n'avons pas estimé fiable l'idée que l'estimateur sarde ait marqué une prédominance avec l'ordre des môts dans son descriptif, ainsi "pâturage et bois" et sera pas différentié de "bois et pâturage").

Les parcelles ne comportant pas de pâturage sont ensuite examinées pour identifier les terres destinées aux cultures (champs, jardins et autres terres cultivées), puis les prés de fauche. Viennent ensuite les espaces boisés, puis les broussailles et bruyères. Enfin, les terrains essentiellement minéraux (roc, roche, rocher, ravines, murgers, glières, etc.) sont regroupés dans la catégorie « Roc ». Une catégorie spécifique est également conservée pour le bâti ainsi que pour les surfaces en eau, marais et rivages. Cette logique hiérarchique est biaisé et orientée avec le parti pris de a) l'importance économique majeur des pâturages pour le Beaufortain en 1730, puis b) les champs permettant des cultures vivrières et c) les prés de fauche de fourrage pour l'hiver, nous avons placé d) les bois avant e) les broussailles (ces dernières ne sont pas du tout improductives car elles permettent de nourrir des caprins et même de fournir une alimentation en hiver pour tout le bétail ) enfin nous avons f) les roc g) les marais... h) le bâti

Cette classification présente l'avantage de restituer les grandes fonctions économiques du territoire tout en restant fidèle à la logique des registres sardes. Elle permet notamment de distinguer les espaces agricoles, les ressources fourragères destinées à l'hivernage du bétail, les pâturages d'altitude, les ressources forestières et les terrains à faible potentiel productif. L'analyse croisée de ces catégories avec l'altitude et le degré de bonté permet ainsi la mise en évidence de l'organisation agropastorale de la communauté de Beaufort au XVIIIe siècle.













\end{document}
